\section{Direct Routes That Did Not Close}

This appendix records the direct routes behind the main class-membership
theorem.  Each entry explains a direct route in four steps: what payment the
route tried to make, why the calculation was worth attempting, where the
mathematics stopped, and what same-solution residue or remaining proof obligation
survived into the Part/Field test.

The order begins with the four-body attempt because that is where the
smoothness-control program first became a concrete circuit:
control the high scales, control energy and enstrophy, use those controls to
produce an exact object, and recover derivative control on that object.
Every later branch below is a refinement of one of those payments.

Each route is written in the same reader-facing order.  First it states the
payment the route tried to make.  Then it shows the calculation or theorem
input that selected the route for testing.  Then it names the mathematical
stopping point.  Finally it
states what survived into the class-membership argument.

\subsection{How to Read the Failure Record}

The reader has already seen the class test.  Part and Field are the
three places where a terminal witness can keep or lose membership.  The
appendix now shows why the positive smoothness program leaves the objects
tested by that grammar.

For that reason every route below has five obligations:
the attempted payment, the calculation or theorem input that selected the route
for testing, the exact
mathematical stopping point, the surviving residue, and the later CM role.
Those obligations are the admission rule for this appendix.  A route belongs
here when it helps the reader identify a terminal object, a member-side readout,
a Part/Field face, or an auxiliary theorem that is still waiting for one of
those jobs.  The appendix serves the main proof by showing where each positive
attempt stopped and what same-solution object or open proof burden it handed
forward.

The order is part of the argument.
The four-body circuit comes first because it is the first complete payment
scheme.
The modern survivor families come next because they split the old scale and
source bills into lower-prefix and upper-tail forms.
The shell, Volterra, signed, selector, source, pressure, Zeno, local-critical,
and endpoint branches then appear in the order in which the unpaid object is
sharpened.
The record closes after the branch table and the CM handoff explain how a
finite-time breakdown claim carries one of these terminal objects into the
Part/Field test.

\subsection{How Repo Source Packets Enter This Appendix}

The live source-frontier, theorem-packet, route-lock, route graph, and proof
assembly surfaces enter this appendix as source rules for failed proof
attempts.  They tell the paper which attempted payments are still current,
which ones are historical, which ones became context only, and which ones
survived as open continuations.  The appendix then writes that source state as a
mathematical sentence on the page: attempted payment, useful calculation,
stopping point, residue, and CM role.

The source-reserve birth packet is the model case.  A zero-radius terminal
residue or selected unpaid infinite donor-refill lands first at
\(\neg Pack_Q\).  A retained positive-scale native donor-square reserve remains
supplier-quarantined until the same ledger pays it, or until the object has
entered the CM test and a first Part or Field Part/Field failure has been
derived.  So source-reserve notes below are context or open continuation until
one of those two payments is written.

The generic Pack-exclusion work order is treated the same way.  It no longer
sets a global demand to prove Pack retention for every terminal obstruction.
Pack work is row-triggered: a named terminal object either certifies
\(\neg Pack_Q\), or it retains Pack and then has to be tested at Part and Field.
The later first-Pack-survival mechanism search records an invention gap: old
source-reserve, no-free-sink, endpoint, material, and bridge labels are failed
attempt or context until a fresh theorem supplies a new input, a falsification
model, and a propagation path to the original-smooth-data Pack question.
The participation-debt reentry supplies the clock language for this audit:
concentration, payment or transfer, and dissipation or return have to be read
inside the same participating fluid.  That clock language is context until it
produces a theorem or a Part/Field landing for a named object.

The current terminal time-face frontier is also recorded here with its exact
limit.  Finite \(L^1\) source mass permits the pulse model
\[
  g_m(s)=m\,1_{(-1/m,0]}(s),
  \qquad g_m(s)\,ds\rightharpoonup\delta_0 .
\]
For any nonzero \(f\in L^2\),
\[
  \int_{-1/m}^0 m e^{(0-s)\nu\Delta}f\,ds\to f
  \quad\text{in }L^2 .
\]
The gold-side continuation would need a temporal thickness theorem, a
super-\(L^1\) source bound, a selected-source tether with absolute continuity,
or a rigid BASAC production theorem that rules out the terminal atom.  When that
gold payment is absent, the appendix records the obstruction and sends any
admitted same-solution residue back to the CM Pack--Part--Field test.

\begin{enumerate}[leftmargin=*,label=\arabic*.]
\item \textbf{Four-body circuit.}
The attempted payment is a loop:
tail control gives \(Q(t)\) control, those controls give compactness, compactness
gives an exact object, and derivative recovery returns to tail control.
The calculation was worth testing because each body is a real mathematical tool:
viscous high-frequency damping, the energy identity, compactness of controlled
approximations, and derivative recovery on a geometric carrier.
The stopping point is circular debt.
Each body needs the hard theorem promised by the next body.
The residue is the first complete obstruction packet: scale transfer, nonlinear
source, exact-object capture, and same-surface derivative drain.
Its CM role is architectural.
It explains why the later proof must test the same terminal object rather than
declare a positive estimate solved.

\item \textbf{Scale barrier.}
The attempted payment is a pointwise nonlinear transfer law strong enough to
make high-frequency derivative tails summable.
The calculation starts from
\(\partial_tE(k,t)=T(k,t)-2\nu\,k^2E(k,t)\) and asks for decay of \(T(k,t)\)
that makes the dominance functional \(\mathcal D_K(t)\) nonnegative on a
high-frequency tail \(K\).
The stopping point is that the energy identity pays
\(\int\|\nabla u\|_2^2\,dt\), not the pointwise cascade law.
The residue is high-frequency nonlinear transfer.
Its CM role is to supply tail, lifted-band, lower-shell, and scale-memory
witness material after a terminal packet has been selected.

\item \textbf{\(Q(t)\) and enstrophy.}
The attempted payment is a one-derivative extension of the kinetic energy law.
The calculation differentiates enstrophy and reaches
\(Y'(t)+\nu Z(t)\le C_0\|\nabla u\|_2^{3/2}\|\Delta u\|_2^{3/2}\), hence
\(Y'(t)\le C_\nu Y(t)^3\) after the Young-margin step spends a fixed fraction of \(\nu Z(t)\).
The stopping point is the vortex-stretching source.
The inequality is local-theory useful and still compatible with finite-time
growth.
The residue is native source debt, later read as source reserve, terminal atom,
or pressure-sustain residue.
Its CM role is to mark source objects as terminal obstructions requiring Pack,
Part, or Field landing.

\item \textbf{Compactness.}
The attempted payment is exact-object production from controlled
approximations.
The calculation is the low/high split: low modes are finite-dimensional once
time compactness is supplied, while high modes are small under a spectral tail.
The stopping point is dependency order.
Compactness consumes tail and enstrophy control; it does not create them.
The residue is same-surface packet capture.
Its CM role is to prevent a terminal object from being treated as proof-bearing
until it is the same Navier--Stokes witness \(Q\).

\item \textbf{Curvature or gradient recovery.}
The attempted payment is derivative drain through geometric structure.
The calculation uses a Riemannian carrier where a positive Ricci term would
control \(|\nabla u|^2\).
The stopping point is surface fidelity.
The Clay problem is periodic and flat, so the Ricci drain is not a free theorem
on the required surface.
The residue is deformation geometry and commutator control.  In CM it remains
context until a same-surface drain theorem or a terminal Part/Field landing is
supplied.

\item \textbf{Modernized four-body packet.}
The attempted payment is the same loop with modern objects:
tail leakage, source reserve, terminal packet capture, and deformation
commutator.
The calculation is the identification of the same survivor through dyadic tail,
velocity tower, and moving-geometry lenses.
The stopping point is the missing contraction or dissipative-margin theorem.
The residue is one shared survivor seen from several directions.
Its CM role is to stop the paper from treating new vocabulary as a new proof.

\item \textbf{Two survivor families.}
The attempted payment is a collapse between the lower-prefix/affine-quotient
family and the upper-tail/gap-kernel family.
The calculation shows one family as scale memory below the active shell and the
other as separated upper-tail leakage.
The stopping point is the missing theorem that either makes one family pay the
other or proves they are the same terminal object.
The residue is a split survivor.
Its CM role is to keep lower-prefix memory and upper-tail leakage separate
until a Part/Field landing identifies their witness role.

\item \textbf{Kernelized lifted band.}
The attempted payment is a gap-kernel gain in the region \(N+M<k<j-4\).
The calculation gets the factor \(2^{-(j-k)}\), then leaves a lower-shell
remainder with an extra \(2^k\) beyond the energy surface.
The stopping point is the unpaid half-derivative.
The residue is a lower-shell gain target.
Its CM role is possible Field or tower readout only after terminal admission.

\item \textbf{Shell gain and half-derivative target.}
The attempted payment is active-window suppression of low-frequency strain.
The calculation normalizes a lower-shell quantity \(A_j(t)\) against
\(\nu2^{2j}\) and asks that the active shell see a small coefficient.
The stopping point is that available estimates are integrated square budgets,
not pointwise active-window bounds.
The residue is a dynamic low-frequency strain theorem.  In CM it feeds retained
tail and Field coherence only after the same terminal packet has been selected.

\item \textbf{Cumulative tail stress and LPAS.}
The attempted payment is a whole-tail version of lower-prefix control.
The calculation packages lower-prefix energy with active-shell square
dissipation in \(P_j^\downarrow2^{-j}D_j^2\).
The stopping point is orientation.
Upper-tail estimates do not pay the lower-prefix active square.
The residue is active-square debt.  In CM it remains tied to source reserve,
scale memory, and same-ledger payment questions.

\item \textbf{Volterra and lower-triangular memory.}
The attempted payment is ordered scale memory.
The calculation turns the lower-prefix object into a Volterra operator and
isolates affine moments \(M_0,M_1\).
The stopping point is a degenerate weight and lack of same-surface coercivity.
The residue is a weighted quotient mode.
Its CM role is terminal-survivor candidate, with exit only after Pack, Part, or
Field failure is derived.

\item \textbf{Signed descent.}
The attempted payment is cancellation before absolute values destroy sign.
The calculation keeps a signed scale derivative and tries to telescope the bad
majorant.
The stopping point is eigendirection alignment: active packets can follow
expanding strain directions.
The residue is missing decorrelation, signed-measure cancellation, or a
nonpositive divergence identity.
Its CM role is return pressure to active-square and source-reserve walls.

\item \textbf{Persistent quotient kill attempts.}
The attempted payment is removal of the affine quotient left by Volterra
memory.
The calculation splits the possible kills into LPAS return, far-corona
control, and affine-moment cancellation.
The stopping point is respectively recycled debt, circular use of the desired
theorem, and a one-derivative shortage.
The residue is affine mass carried by boundary, spill, or edge terms.
Its CM role is terminal quotient evidence requiring a first-face certificate.

\item \textbf{Projected flow, TPS, and SG.4.}
The attempted payment is a parabolic or two-point defect equation with a
coercive energy identity.
The calculation gives a desired estimate
\(|N_1|+|N_2|\le\epsilon\,\mathrm{dissipation}+C_\epsilon\Psi(E_D)\).
The stopping point is selector production.
Norm control can lose sign, direction, pair membership, and one-sided growth.
The residue is selector-good lower-envelope debt.  In CM it remains context
unless it produces the selected packet or an exact one-sided theorem.

\item \textbf{Source-wall deletion cycle.}
The attempted payment is direct deletion of the nonlinear source from the
old \(Q(t)\) obstruction.
The calculation repeatedly charges source pulses to native ledgers.
The stopping point is that each charge exposes a sharper source object.
The residue is terminal native source concentration.
Its CM role is to admit source objects to the witness test rather than hide
them behind a missing forward estimate.

\item \textbf{BASAC and terminal atom.}
The attempted payment is bounded source mass strong enough to prevent terminal
time atoms.
The exact test sequence \(g_m=m1_{(-1/m,0]}\) satisfies
\(\|g_m\|_{L^1}=1\), \(g_m(s)\,ds\rightharpoonup\delta_0\), and
\(\|g_m\|_{L^p}=m^{1-1/p}\to\infty\) for every \(p>1\).
The stopping point is lack of uniform integrability or anti-atom control.
The residue is a zero-thickness native source atom.  In CM it becomes
zero-radius Part/Field face evidence after CM-test admission.

\item \textbf{Pressure sustain and pure pressure residue.}
The attempted payment is reduction of pressure debt to native source debt.
The calculation separates source-carried pressure from a singular pressure
part orthogonal to native source.
The stopping point is that Calderon--Zygmund pressure is spatial and tensorial;
the exact missing input is a same-fluid pressure-source ancestry theorem for
the selected pressure contribution.
The residue is pure pressure sustain.
Its CM role is Field diagnostic or pressure-side support after pressure exits
are isolated.

\item \textbf{Angular saturation.}
The attempted payment is cancellation of pressure Hessian lobes by angular
partners.
The calculation sees that tensor trace cancellation can coexist with a selected
positive terminal lobe.
The stopping point is partner legality: negative partners may diffuse, fall
below threshold, leave the family, or miss the selected weights.
The residue is pressure-lobe debt.  In CM it remains pressure-tethering or pure
pressure Field-failure context until a same-witness is derived.

\item \textbf{Transported boundary and no incoming flux.}
The attempted payment is elimination of source birth through same-fluid
transported cylinders.
The calculation removes artificial fixed-frame boundary inflow.
The stopping point is interior time-face concentration.
An atom can be born inside the shrinking transported cylinder.
The residue is no-incoming-flux plus interior-production debt.
Its CM role is Part or Field support depending on retained carrier and
flux role.

\item \textbf{Quantized parent and Zeno ancestry.}
The attempted payment is a monotone resource spent by each same-fluid refill.
The calculation tries radius, heat-time depth, energy, dissipation, labels, and
donor balance.
The stopping point is infinite divisibility down the terminal scale tree.
The residue is terminal Zeno ancestry.
Its CM role is carrier-window support zero-radius boundary unless a retained positive-scale
carrier is separately produced.

\item \textbf{Local energy and elliptic time smearing.}
The attempted payment is exclusion of terminal source atoms by local energy or
pressure ellipticity.
The calculation shows local energy gives signed terminal measures, while
spatial elliptic operators act time-slice by time-slice.
The stopping point is that time concentration can survive both tools.
The residue is trace, flux, or anti-concentration debt.
Its CM role is source-wall support.

\item \textbf{All reduced source-wall routes.}
The attempted payment is a finite list of supplier theorems:
temporal source integrability, positive-source Carleson control, parent
quantization, pressure-source tethering, pure-pressure Liouville, and terminal
pressure-strain legality.
The stopping point is that none is installed unconditionally from the original
data.
The residue is the reduced source-wall obstruction list.
Its CM role is honesty: these are supplier doors, not completed proof payments.

\item \textbf{Source reserve birth.}
The attempted payment is charging first-created donor-square reserve to legal,
residual, or inherited ledgers.
The calculation chooses a first terminal same-fluid window \(W\) and contracts
past reserve by minimality.
The stopping point is the first-created native reserve left on the retained
ledger.
The residue is positive-scale donor-square reserve.
Its CM role is supplier-quarantined until a same-ledger payment theorem or
CM-test Part/Field failure is supplied.

\item \textbf{Zero-moment signed pair.}
The attempted payment is cancellation of reserve by zero first moment.
The calculation splits \(J=J_+-J_-\) with equal first moments and fixed
positive mass: there is \(c_0>0\) such that
\(|J_+|(W),|J_-|(W)\ge c_0\) on the retained window \(W\).
The stopping point is that two packets can exchange through the same retained
ledger while moments cancel on heat subwindows.
The residue is a same-ledger signed-pair no-free-sink theorem.
Its CM role is retained reserve support, not exit without a Part/Field landing.

\item \textbf{Public \(L^3\), Duhamel, and native source.}
The attempted payment is a local critical translator, not fresh deletion.
The calculation splits Duhamel into heat ancestry and nonlinear native-source
response on the same witness ledger.
The stopping point is not the local translator; it is whether the terminal
mass is already on the same ledger.
The residue is terminal Duhamel response mass.
Its CM role is installed local witness translation into Pack, Part, legal, or
Field-facing readout.

\item \textbf{Terminal Leray saturation.}
The attempted payment is conversion of local retained source work into a
saturated square.
The calculation isolates the signed-partner defect after writing
\(\widetilde w_j=w_j+\delta_j\).
The local \(L^3\) translator handles the same-family signed defect.
The stopping point is the independent pressure-Leray route, which still needs
a source or donor edge theorem.
The residue is pressure-lobe tether debt.
Its CM role is closed only inside the local \(L^3\) family.

\item \textbf{Critical \(H^{1/2}\).}
The attempted payment is a frequency-sensitive critical translator.
The calculation separates amplitude, which routes toward \(L^3\), from pure
oscillation, which is Field-facing only after the same-ledger and
oscillation-to-Field lemmas are supplied.
The stopping point is missing same-ledger extraction and oscillation-to-Field
lemmas.
The residue is conditional \(H^{1/2}\) field defect.
Its CM role is conditional translator only.

\item \textbf{Synthesis and branch table.}
The attempted payment at the end is reader custody.
The paper must not leave the branches as a pile of names.
The common-pattern section identifies the recurring motion: remove legal
pieces, isolate a terminal residue, sharpen the survivor, then test it.
The branch table records each route's forward residue and later CM role.
The closure states the actual handoff: a finite-time breakdown claim must
carry one of these terminal objects, and the main proof tests that object
through Part and Field.
\end{enumerate}

\subsection{Expanded Branch Work Packets}

The following packets are the reader-facing expansion of the contract above.
Each packet says what the route was trying to buy for a positive smoothness
proof, what calculation made that purchase look available, why the purchase
did not go through, and what object the CM proof inherits.

\subsubsection{Four-body circuit.}

The four-body circuit was the first complete attempt to make the positive
smoothness route pay its own bills.
It had four bodies.
The first body was tail control.
The second body was a controlling quantity \(Q(t)\) built from energy and
enstrophy.
The third body was the compactness statement: controlled approximations, with
the required topology and pressure controls, select an exact Navier--Stokes
object.
The fourth body was derivative recovery on that exact object.

The intended motion was the following four-statement implication chain.
Tail control supplies high-frequency derivative-mass control.
The \(Q(t)\) theorem binds that control to a finite energy-enstrophy quantity on
the time interval under study.
The compactness theorem selects an exact terminal object.
The derivative-recovery theorem recovers the bound needed to restart the
high-frequency tail estimate.
If all four statements are proved on the periodic surface, the classical
continuation criterion closes the proof.

The first body starts from the spectral energy balance.  Write \(E(k,t)\) for
the energy density at frequency \(k\), and write \(T(k,t)\) for the nonlinear
transfer into that shell.  The model identity has the form
\[
  \partial_t E(k,t)=T(k,t)-2\nu\,k^2E(k,t).
\]
The viscous term has the right sign and the right frequency weight.
A direct proof has to state this as a tail inequality on the same frequencies it
later consumes.  With \(T_+(k,t)=\max(T(k,t),0)\), and for a cutoff \(K\), set
\[
  Y_{>K}(t):=\int_K^\infty k^2E(k,t)\,dk
\]
and
\[
  \mathcal D_K(t)
  :=
  2\nu\int_K^\infty k^4E(k,t)\,dk
  -
  \int_K^\infty k^2T_+(k,t)\,dk .
\]
The scale payment is not a pointwise profile for \(T\) and not an active-scale
picture.  It is a proof that \(\mathcal D_K(t)\ge0\), or that
\((\mathcal D_K(t))_-\) is integrable with the required constants, on the
selected tail.

The second body asks the same question in physical-space energy language.
Let
\[
  Y(t)=\frac12\|\nabla u(t)\|_2^2 .
\]
Testing the equation against \(-\Delta u\), integrating by parts, and using
\(\nabla\cdot u=0\) gives the enstrophy identity
\[
  Y'(t)+\nu\|\Delta u(t)\|_2^2
    = \int_{\mathbb T^3}(u\cdot\nabla)u\cdot\Delta u\,dx .
\]
The left side is the desired derivative payment.
The right side is the source that has to be paid.
H\"older, Sobolev interpolation, and Young bound the source term by the
conditional local-continuation inequality
\[
  Y'(t)\le C_\nu Y(t)^3 .
\]
That inequality explains local persistence, and it also explains the global
gap.
The comparison equation can blow up in finite time.
The \(Q(t)\) body must do more than repeat enstrophy control; it has to find a
new payment for the source term.

The third body is compactness.
If the tail and source payments were already available, compactness would still
need a specified extraction topology and a same-solution attachment theorem.
One splits an approximating sequence into low and high modes.
Low modes are finite-dimensional after a time-compactness bound is supplied.
High modes are small only after a uniform selected-tail estimate is proved.
Then strong convergence passes the quadratic term,
\[
  u_n\otimes u_n\longrightarrow u\otimes u,
\]
and the limit remains a solution of the same equation only after the weak
equation, pressure normalization, and initial-data/uniqueness attachment are
kept on the original Navier--Stokes surface.
This body is essential because the Clay problem is about the original
Navier--Stokes solution, not about an auxiliary object that merely resembles
it.
The compactness payment is so exact-object capture.

The fourth body is derivative recovery.
Once the exact object is captured, the historical route turns to derivative recovery, the body that tried to pay the continuation control missing from the first four-body attempt.
In the early geometric version of the route, the missing drain theorem was a
positive same-surface curvature or deformation term in a gradient evolution law,
\[
  \partial_t|\nabla u|^2
    = -\mathcal D(\nabla u) + \nu\Delta |\nabla u|^2
      + \mathcal N(\nabla u),
\]
where \(\mathcal D\) would be a same-surface drain and \(\mathcal N\) would be a
controlled nonlinear remainder.
The periodic problem supplies viscosity and pressure coupling.
The exact missing input is a same-surface derivative recovery theorem: a
coercive drain with nonlinear and pressure terms controlled on the same
surface.

The circuit fails because each body spends a theorem that another body was
supposed to earn.
Tail control needs the nonlinear source to be paid.
The \(Q(t)\) body needs a source estimate stronger than the local enstrophy
comparison.
Compactness needs tail and source control before it can pass the quadratic
term cleanly.
Derivative recovery needs the exact object and a same-surface drain.
No body is meaningless, and no body is a completed global proof.  The attempted
loop reveals one obstruction packet: tail payment, source payment, exact-object
capture, and same-surface derivative recovery.

\paragraph{Four-body debt lemma.}
Suppose a positive smoothness proof uses the four-body circuit above.
Then the proof has to provide four distinct payments on the same periodic
Navier--Stokes branch: a summable high-frequency tail estimate, a source
payment strong enough to prevent the enstrophy comparison from becoming the
only available control, compactness that captures the same solution object, and
derivative recovery on that same object.
If one payment is absent, the circuit leaves a terminal residue instead of a
global continuation theorem.

\paragraph{Proof.}
The tail estimate is needed before the high modes can be treated as small.
The source payment is needed before the derivative energy can be continued by
an ordinary differential comparison.
The compactness step is needed before the limiting object can be charged as the
same Navier--Stokes branch.
The recovery step is needed before the argument can return from the captured
object to the derivative bounds required by classical continuation.
Removing any one of these payments breaks the loop at the point where that
payment is spent.
The remaining object is then a residue of the failed positive route.

\paragraph{CM handoff.}
The class-membership proof inherits the four-body failure as a discipline for
terminal witnesses.
A terminal breakdown claim cannot be accepted merely because it names tail
growth, source concentration, a compactness limit, or a geometric drain.
It has to present the same Navier--Stokes object \(Q\).
The paper then asks the Field-certification questions.
Does the object still have a positive packet?
Does the same pressure-viscosity equation still act on that packet?
Does the branch still have a coherent field on a positive scale?
The old four-body circuit explains why those are the right questions.
They are the four old payments rewritten as membership tests on the selected
terminal object.

\paragraph{Scale barrier.}
The scale route tried to turn the sentence ``viscosity wins at small scales''
into a summable derivative estimate.
The useful calculation starts from the spectral balance
\[
  \partial_tE(k,t)=T(k,t)-2\nu\,k^2E(k,t)
\]
on the zero-force branch.
If the nonlinear transfer \(T(k,t)\) decayed fast enough, then above a critical
scale the viscous term would dominate and the derivative tail
\(\int k^2E(k,t)\,dk\) would stay finite.
That is a legitimate proof dream because it identifies an exact payment:
high-frequency transfer must be weaker than viscous loss.
The failure is that the energy law gives an integrated dissipation budget, not
a pointwise cascade law with the needed exponent.
The residue is a precise high-frequency nonlinear transfer packet.
In the later paper this packet reappears as lifted-band leakage, lower-shell
gain, cumulative tail stress, and scale memory.
It enters CM only after a selected terminal object carries the tail record and
then loses Part or Field.

\paragraph{\(Q(t)\) and enstrophy.}
The \(Q(t)\) route tried to make kinetic energy carry one more derivative.
The kinetic half is exact and harmless, but differentiating enstrophy exposes
the source:
\[
  \frac{dY}{dt}+\nu Z=\int_{\mathbb T^3}(u\cdot\nabla)u\cdot\Delta u\,dx .
\]
Sobolev and interpolation estimate the right side by a power of
\(\|\nabla u\|_2\) and \(\|\Delta u\|_2\), after which the Young-margin step leaves
\[
  \frac{dY}{dt}\le C_\nu Y^3 .
\]
That inequality is enough for local theory and still permits finite-time
growth.
So the route does not fail because the calculation is sloppy.
It fails because the nonlinear vortex-stretching source is exactly where the
global theorem would have to be paid.
The old missing source later returns as terminal native source, source reserve,
BASAC atom, pressure-sustain residue, and Duhamel response.
The CM use is disciplined: a source object is not an exit because it is
dangerous; it is proof material only when the selected terminal witness loses a
class face.

\paragraph{Compactness.}
The compactness body tried to turn prelimit control into an exact solution
object.
The route is mathematically honest.
Split an approximating sequence into low and high frequencies.
The low-frequency part is finite-dimensional after time compactness is
supplied, while the high-frequency part is small if the spectral tail is small.
Strong local \(L^2\) convergence then passes the quadratic term
\[
  u_n\otimes u_n\to u\otimes u .
\]
This is why compactness belongs in the record: it tells the proof how a
controlled packet becomes the same Navier--Stokes object rather than a loose
limit.
The failure is that compactness is downstream.
It spends the scale tail and enstrophy controls; it cannot manufacture them.
When the PDF later speaks of packet capture, the reader should hear this
lesson.
Capture is useful only after the packet is already on the correct surface.
In CM, this becomes the rule that auxiliary estimates do not enter the proof
until they identify the same terminal witness \(Q\).

\paragraph{Curvature or gradient recovery.}
The curvature route tried to recover derivative control from geometry.
On a curved carrier, the model gradient inequality used by the branch contains
a Ricci term: there is a geometry-dependent constant \(C_g\) for which
\[
  \partial_t|\nabla u|^2
  \le
  -\mathrm{Ric}(g)(\nabla u,\nabla u)
  +\nu\Delta_g|\nabla u|^2
  +C_g|\nabla u|^3 .
\]
A positive Ricci contribution would be exactly the missing derivative drain.
That is why this branch looked powerful.
It was trying to pay the fourth body of the loop and return control to the
scale barrier.
The failure is not cosmetic.
The Clay problem here is periodic and flat.
A positive-Ricci theorem on a different carrier is not the same-surface
periodic theorem the proof needs.
The survivor is a demand for a replacement theorem inside the actual
Navier--Stokes geometry: deformation-gradient control, commutator drain, or
another same-surface mechanism.
Until that theorem exists, the curvature branch is support and motivation, not
a proof payment.

\paragraph{Modernized four-body packet.}
The modern packet keeps the four old jobs and replaces their vague names by
the objects that actually survived.
Scale control becomes tail leakage.
The \(Q(t)\) source becomes source reserve.
Compactness becomes terminal packet capture.
Curvature becomes deformation geometry and commutator control.
The progress is real because it stops treating each new vocabulary item as a
new proof.
The same obstruction is being read through several instruments: dyadic tail,
velocity tower, moving coordinates, and terminal source account.
The failure is also clearer.
The packet still needs a same-surface contraction or dissipative-margin law strong
enough to make the survivor pay for itself.
Without that law, the packet is a sharper description of debt, not a discharge
of debt.
The CM proof uses the vocabulary here to keep this survivor in its correct
role.
If a later object is merely a modern name for the same unpaid survivor, the
paper must carry it as residue until Part or Field makes the proof
decision.

\paragraph{Two survivor families.}
After the modern four-body pass, the residue splits into two large families.
The lower-prefix family reads the obstruction as memory from scales below the
active shell: affine quotient modes, Volterra primitives, lower-shell strain,
and scale-memory terms.
The upper-tail family reads it as leakage from separated shells above or around
the active band: gap kernels, lifted-band packets, upper boundary effects, and
kernelized tail dissipation.
The key point is that both families descend from the same old scale bill, but
they do not currently pay one another.
The missing theorem is one of three exact alternatives: lower-prefix memory
creates a coercive drain; upper-tail leakage is blocked by a kernelized barrier;
or the two descriptions are one terminal survivor in different coordinates.
Until then, the paper must keep them separate.
In CM terms, neither family is an exit without a selected terminal packet and a
first Part/Field failure.
Each remains a witness candidate until that test is paid.

\paragraph{Kernelized lifted band.}
The lifted-band route attacks a precise shell region, \(N+M<k<j-4\), where
off-diagonal interactions have a gap.
The calculation gains a real kernel factor \(2^{-(j-k)}\).
That gain matters because it says separated shells do not interact with full
strength.
After the Young-margin split, however, the lower shell carries a remainder with an
extra \(2^k\) beyond the available energy surface.
The route so gets closer to the desired estimate but stops at an unpaid
half derivative.
This is the kind of failure the PDF must preserve carefully.
The branch is not worthless, and it is not complete.
It gives an exact leftover object: an intermediate lower-shell gain target.
In the CM proof, that object can support Field or tower readout only after it
belongs to a selected terminal witness.
Before that admission, it is a forward-positive survivor.

\paragraph{Shell gain and half-derivative target.}
The shell-gain branch tries to pay the extra lower-shell factor left by the
lifted-band calculation.
It normalizes a low-frequency strain quantity against the viscous scale
\(\nu2^{2j}\) and asks the active high shell to see a small coefficient.
The calculation has the right scaling: the high shell dissipates on a
\(2^{-2j}\) time scale, so lower-shell strain has to be measured against the
same parabolic clock.
The stopping point is the difference between an integrated square budget and a
pointwise active-window bound.
The installed estimates say that strain is affordable on average.
They do not say that every active window satisfies the coefficient bound
\(B_j(t)\le\varepsilon\) needed to keep growth inside the viscous margin.
The survivor is dynamic low-frequency strain suppression.
In the PDF this branch should teach the reader why the active-window coefficient
bound is a theorem target, not a rhetorical strengthening of viscosity.

\paragraph{Cumulative tail stress and LPAS.}
Cumulative tail stress tries to stop chasing one shell at a time.
It packages the lower prefix and active shell into a whole-tail object, whose
dyadic shadow is the lower-prefix active square \(P_j^\downarrow2^{-j}D_j^2\).
The exact advantage is that this object records the whole lower-scale history
that one-shell estimates miss.
The failure is orientation.
Many available estimates are upper-tail estimates; they control what happens
above the active index.
LPAS asks the opposite question: how much lower-prefix energy is multiplying
the active square at the bad shell.
A crude bound by total dissipation throws away the active-shell structure and
does not close the dissipative-margin estimate.
The residue is active-square debt.
It later reappears as source reserve and lower-memory pressure.
In CM, LPAS remains context unless the same terminal packet carries the active
square and loses one of the three class requirements.

\paragraph{Volterra and lower-triangular memory.}
The Volterra route recognizes that lower-prefix effects are ordered in scale.
They are not arbitrary bilinear noise.
In logarithmic scale the lower prefix becomes a triangular memory operator:
earlier or lower scales feed the active shell in one direction.
That observation lets the proof isolate affine moments such as \(M_0\) and
\(M_1\).
This is a real reduction.
A broad lower-prefix obstruction becomes a small quotient mode.
The exact missing input is a weighted quotient control theorem.
The weight can degenerate, and the installed estimates do not force the affine
moments to vanish on the same surface.
The survivor is a weighted quotient mode carried by the terminal branch.
The CM proof uses this as a warning against cheap bad-object language.
An affine quotient is not class exit until the terminal object has been tested
and the first lost face has been identified.

\paragraph{Signed descent.}
Signed descent tries to preserve cancellation before the proof takes absolute
values.
The calculation keeps a signed scale derivative visible and tries to telescope
the dangerous term, so that a positive majorant never has to be paid.
This is the right instinct: many failed estimates become impossible only after
the proof has destroyed sign too early.
The stopping point is the strain geometry.
The leading form has no fixed sign, and active wave packets can align with
expanding eigendirections of the low shell.
One-sided weights do not by themselves prevent that alignment.
The survivor is so a missing decorrelation theorem, a signed-measure
cancellation theorem, or a nonpositive divergence identity.
When those are absent, the proof returns to the active-square/source-reserve
wall.
The PDF must show this return so the reader does not mistake signed notation
for signed control.

\paragraph{Persistent quotient kill attempts.}
The quotient-kill attempts ask whether the affine mode left by Volterra memory
can be destroyed directly.
The first route sends the problem back to CTS/LPAS, which means it has not
escaped the lower-prefix active-square wall.
The second route invokes a far-corona theorem strongly enough to be circular:
it uses the kind of tail control the branch was supposed to produce.
The third route isolates affine moments and tries to spend a signed derivative,
but it remains one derivative short.
The important point for the reader is that these are not three independent
failures.
They are three tests of the same survivor.
Boundary, spill, and edge terms all carry versions of the affine mass.
The CM role is terminal quotient evidence after witness admission.
Before that admission, the quotient is proof pressure, not a completed exit.

\paragraph{Projected flow, TPS, and SG.4.}
The projected-flow and TPS branch tries to put the defect into an equation
where parabolic coercivity or two-point geometry can see it.
The desired projected-flow estimate has the form
\[
  |N_1|+|N_2|\le \epsilon\,\mathrm{dissipation}+C_\epsilon\Psi(E_D).
\]
The two-point strain side has a parallel theorem target: selected pair geometry
must produce the good lower envelope needed by SG.4.
The failure is selector production.
A norm can bound the size of a defect while losing exactly the information the
selector needs: sign, direction, membership in the selected pair family, and
one-sided growth.
That is why SG.4 cannot be promoted to route authority merely because it names
a sharp downstream obstruction.  In CM it remains readout pressure until it
lands the exact same Navier--Stokes packet in Part or Field.

\paragraph{Source-wall deletion cycle.}
The source-wall deletion cycle is the old \(Q(t)\) obstruction made terminal.
Each phase tries to delete or pay the nonlinear source.
The recurring pattern is stable: find a source pulse, charge it to a native
ledger, remove legal or inherited pieces, and discover a sharper survivor.
This is why the PDF cannot present source work as a single failed estimate.
The source wall is a machine for producing terminal objects.
Sometimes the object is a reserve.
Sometimes it is a zero-thickness atom.
Sometimes pressure separates off a residue that is not native source.
The stopping point is that forward deletion keeps reaching one more terminal
source object rather than a global positive bound.
The CM route uses this by admitting the terminal source object to the witness
test.
It asks whether the object still has Part and Field, not whether the old
positive deletion theorem has magically been proved.

\paragraph{BASAC and terminal atom.}
BASAC tries to control source mass at the terminal layer.
The model density \(g_m=m1_{(-1/m,0]}\) explains the obstruction in one line:
the \(L^1\) mass stays bounded while
\[
  \|g_m\|_{L^p}=m^{1-1/p}
\]
for every \(p>1\), so every super-\(L^1\) norm diverges.
So finite source mass can converge to a terminal time atom.
The branch fails at uniform integrability, reverse Holder residence, Morrey
anti-atom control, or a Liouville mechanism strong enough to rule out the
atom.
The survivor is a zero-thickness native source residue.
This is exactly the kind of object that makes honesty essential.
Calling it impossible would import the missing anti-atom theorem.
Calling it an exit at that point would skip the CM Part/Field test.
The paper must instead say what the object is, why the positive deletion route
stopped there, and how CM tests it as zero-radius Part/Field face evidence after
terminal admission.

\paragraph{Pressure sustain and pure pressure residue.}
The pressure branch asks whether pressure debt can be reduced to native source
debt.
The pressure equation is elliptic in space and tied to quadratic gradients, so
the required bridge is a pressure-source tether, not a source-shadow reading.
The failure is same-fluid ancestry.
Calderon--Zygmund structure acts on spatial slices and tensor components; it
requires a pressure-source tether before the pressure lobe belongs to the same
native source current.
The branch so splits pressure into a source-carried part and a pure
pressure residue orthogonal to native source.
The missing theorem is either a pressure-source tether or a pure-pressure
Liouville exclusion.
Until one exists, pressure sustain is a diagnostic survivor.
Its CM role is usually Field-side support after pressure exits are isolated,
because the issue is coherent readout and participation of the pressure lobe
inside the selected same-solution packet.

\paragraph{Angular saturation.}
Angular saturation tries to use cancellation in the pressure Hessian.
The trace can cancel over directions, and tensor structure can balance lobes
globally.
The selected terminal ledger, however, may keep one positive lobe with its own
weights, cutoffs, and family membership.
The negative partners needed for cancellation may diffuse, fall below
threshold, sit outside the selected family, or arrive with incompatible
weights.
That is the stopping point.
Global tensor cancellation pays the selected positive bill only after partner
legality is proved for the selected weights, cutoffs, and family membership.
The survivor is pressure-lobe debt together with a partner-legality theorem
stack: partner saturation, pair-weight defect charge, and collar legality.
In CM language this branch is not a new primitive face.
It is pressure-side support for either a source tether, pure pressure
Liouville, or a Field failure once the selected pressure object is held on the
same witness record.

\paragraph{Transported boundary and no incoming flux.}
Transported cylinders are the right geometry for same-fluid ancestry.
They remove artificial fixed-frame boundary inflow, so the proof no longer
charges source birth to a coordinate artifact.
The required theorem says that incoming positive source flux through the
transported boundary tends to zero.
The failure is that boundary control does not prevent interior time-face
concentration.
A source atom can be born inside the transported cylinder as the terminal layer
shrinks.
The local energy balance then still has a terminal source term to pay.
The missing theorem is a combined no-incoming-flux and interior-production
exclusion statement.
The CM role depends on what survives.
No positive carrier points to Pack.
Broken same-equation participation points to Part.
A retained carrier with no coherent positive-scale readout points to Field.
The PDF needs this split because no-flux alone is not the proof.

\paragraph{Quantized parent and Zeno ancestry.}
The Zeno route tries to make infinite same-fluid refill spend a finite resource.
If each refill edge costs at least \(\eta>0\), infinitely many refills would
contradict finite original data.
The proof tries several resources: radius, heat-time depth, energy,
dissipation, same-fluid labels, and donor balance.
Each one fails for a precise reason.
Radii can shrink.
Heat times can be summable.
Energy and dissipation can split over infinitely many descendants.
Labels preserve ancestry without creating a decreasing finite alphabet.
Donor balance telescopes over finite truncations and leaves terminal leaves.
The survivor is terminal zero-radius Zeno ancestry.
In CM, Zeno is not a fourth face.
A pure zero-radius endpoint is carrier-window support boundary unless a positive-scale
carrier is separately retained; retained variants must still be routed through
Part or Field.

\paragraph{Local energy and elliptic time smearing.}
The local-energy branch asks whether terminal source atoms can be excluded by
the local energy inequality.
The inequality gives compactness and signed terminal defect measures, but it
does not by itself remove the selected positive native source component.
The elliptic branch asks whether pressure or Calderon--Zygmund spreading smears
the atom.
The obstruction is that spatial elliptic operators act one time slice at a
time.
A density \(m1_{(-1/m,0]}(s)\rho(y)\) remains concentrated in time after a
spatial singular integral changes \(\rho(y)\).
Thus local energy and elliptic pressure are real auxiliary tools, but neither
one supplies the missing time anti-concentration theorem.
The residue is trace, flux, or temporal spreading debt.
The CM proof uses this as source-wall evidence, not as a finished deletion of
terminal atoms.

\paragraph{All reduced source-wall routes.}
The all-routes pass is an honesty inventory.
It names the finite set of supplier doors still available after source-wall
reduction: uniform temporal source integrability, positive source Carleson or
remainder depletion, parent quantization or minimal ancestry compactness,
pressure-source tethering, pure-pressure Liouville, and terminal pressure-strain
legality.
The important content is negative as well as positive.
None of these is installed unconditionally from the original data inside this
paper.
so none may be silently used as if it has paid the source wall.
The residue is the shared terminal zero-thickness source concentration plus
pressure-side variants.
The CM role is to quarantine supplier claims until they either prove a real
positive theorem or attach a selected terminal object to Part or Field.
The honesty burden is highest here because these survivors are tempting to
rename as progress before they have paid a face.

\paragraph{Source reserve birth.}
The source-reserve branch chooses the first terminal same-fluid window where
donor-square reserve is born.
Minimality is supposed to contract inherited past reserve, while legal and
child residual terms are routed to existing ledgers.
The calculation leaves a first-created native reserve on the retained ledger.
That survivor is sharper than the old source pulse because it is not merely
mass; it is positive-scale donor-square reserve created on the same terminal
window.
The direct proof stops because the installed ledgers do not pay this first
creation.
The branch so cannot be called deleted.
It is supplier-quarantined until a same-ledger payment theorem is proved.
Alternatively, once admitted to the CM test, it can support a first Part/Field
failure.
The PDF has to preserve that fork: payment theorem on the positive side,
Part/Field failure on the contrapositive side.

\paragraph{Zero-moment signed pair.}
The zero-moment branch shows why signed cancellation is not the same as
reserve disappearance.
One can split a current as \(J=J_+-J_-\), arrange equal first moments, and
still retain positive and negative packets on the same ledger with
\(|J_+|(W),|J_-|(W)\ge c_0\) for a fixed \(c_0>0\).
The square reserve sees mass even when the signed moment cancels.
Most easy cases route away: wrong carrier, shrinking terminal windows, or
broken participation.
The hard case is same retained ledger, positive scale, reserve bounded below by
a fixed \(c_0>0\), and
zero first moment on heat subwindows.
The missing theorem is a same-ledger signed-pair no-free-sink theorem or a
two-tower donor depletion theorem.
Until then, zero moment is not deletion.
In CM it is retained reserve support and can become proof-bearing only after a
first Part or Field Part/Field failure is derived.

\paragraph{Public \(L^3\), Duhamel, and native source.}
The \(L^3\) branch is different from the source deletion branches because it
is a translator.
The public critical criterion gives terminal critical mass.
Duhamel splits the localized solution into heat ancestry and nonlinear
response:
the heat side checks where the packet came from, and the nonlinear side is
localized by adjoint heat/Leray response into the native-source argument.
The stopping point is not that the local translator fails.
The condition is that the terminal \(L^3\) mass must already live on the same
witness ledger.
Once it does, the installed local translator routes the mass into Pack, Part,
legal material, or Field-facing readout.
This branch is a CM-facing translation rule for a selected terminal concentration.
The PDF should make that role explicit so the reader does not confuse public
critical theory with the source-wall deletion program.

\paragraph{Terminal Leray saturation.}
Terminal Leray saturation sits beside the \(L^3\) translator.
After legal commutator removal, the retained local source work is compared to a
saturated square.
Writing \(\widetilde w_j=w_j+\delta_j\) isolates a signed-partner defect.
Inside the local \(L^3\) family, the signed-partner no-free-defect theorem pays
that defect, so the local translator can close its own branch.
The independent pressure-Leray route remains open.
A pressure lobe still needs a transported source or donor-edge theorem before
it can be claimed as native source work.
The survivor is pressure-lobe tether debt.
The CM role is narrow: terminal Leray saturation closes what belongs to the
local \(L^3\) translator, while pressure-Leray material remains support until
it lands a selected same-solution packet or proves its tether.

\paragraph{Critical \(H^{1/2}\).}
The \(H^{1/2}\) branch is a conditional critical translator.
It sees both size and frequency placement, so it can split terminal behavior
more finely than \(L^3\).
The amplitude reading routes toward the installed local \(L^3\) translator.
The pure oscillatory reading is different: amplitude remains controlled while
frequency organization destroys coherent readout.
With Pack and Part retained, that pure oscillatory branch has the shape of a
Field failure, but it is not a proved Field failure yet.
The stopping point is that the necessary lemmas are not installed:
same-ledger \(H^{1/2}\) extraction, oscillation-to-Field exit, and amplitude to
\(L^3\) translation.
The residue is a conditional \(H^{1/2}\) field defect.
The PDF must mark it conditional.
It can guide the reader toward a possible Field landing, but it cannot be used
as completed class exit.

\paragraph{Synthesis.}
The final synthesis is not a conclusion by mood.
It is the place where the reader should see the common form of the whole
failure record.
Each positive branch removes the local or legal pieces it can remove, isolates
an edge, collar, source, reserve, quotient, pressure, Zeno, or critical packet,
and then discovers that the smaller packet still needs a payment theorem.
The CM proof does not deny those packets.
It changes the question asked of them.
Does the selected terminal object still carry a positive same-solution packet?
Does the Navier--Stokes law still participate on that packet?
Does the resulting field have a coherent positive-scale readout?
The branch table and closure must leave the reader with that exact handoff:
the failure history identifies what a finite-time breakdown claim would have to
carry; the main proof tests that carried object through Part and Field.

\subsection{How The Record Enters The Proof}

The failure record is not an appendix of anecdotes.
It is the back half of the reader's route through the proof.
The main text has already stated the three requirements that define the working
class: a positive material packet, same-equation participation on that packet,
and coherent positive-scale field readout.
Those requirements are abstract until the reader sees what a terminal
breakdown attempt would actually carry.
The branch record supplies that missing concreteness.
It shows the reader the objects produced by honest positive proof attempts:
tail packets, source atoms, affine quotient modes, pressure lobes, Zeno
ancestry, signed-pair reserve, critical Duhamel response, and conditional
\(H^{1/2}\) oscillation.
The paper then asks the same question of each object.
Does this object still give the same Navier--Stokes branch, or has it lost the
structure needed to count as that branch?
That is why the record cannot be replaced by a short list of route names.
The reader has to see the mathematical birth of each object before the CM test
can feel like a proof rather than a classification trick.

The first phase is the four-body root.
This phase teaches the reader that the original positive proof did not fail
because it was directionless.
It had a coherent payment scheme: scale suppression, enstrophy control,
compact exact-object production, and derivative recovery.
The page has to show that the scheme failed at actual mathematical joints.
Scale suppression needs a nonlinear transfer law.
Enstrophy control needs vortex-stretching source payment.
Compactness needs the previous controls already in hand.
Derivative recovery needs a same-surface drain, not a geometric theorem on a
different carrier.
Once this is visible, the later proof has a different shape.
It no longer searches for a single magical estimate that overwrites the four
debts.
It asks what terminal object is left when those debts cannot be paid directly.
The four-body root so belongs before the branch families.
It gives the reader the grammar for recognizing every later survivor as one of
the old unpaid payments in a sharper form.

The second phase is the scale and tail phase.
The lifted-band, shell-gain, cumulative-tail, LPAS, and upper-tail routes all
come from the scale barrier.
They should not be scattered as unrelated lemmas.
They are attempts to locate where the cascade payment fails.
The lifted-band route finds a real gap-kernel gain and then leaves an unpaid
lower-shell factor.
The shell-gain route tries to suppress that factor by comparing low-frequency
strain to the parabolic clock of the active shell.
The cumulative-tail and LPAS route then realizes that the problem is not one
shell but a lower-prefix active square.
The reader needs to see this progression because it prevents two errors.
First, it prevents the proof from calling every sharper estimate a new route.
Second, it prevents the proof from throwing the scale survivor away as generic
failure.
The scale phase produces a terminal tail or lower-prefix object that may later
be read through Field coherence, but only after the object is selected on the
same witness ledger.

The third phase is lower-triangular memory.
Volterra, affine quotient, signed descent, and quotient-kill attempts are the
memory-side continuation of the scale phase.
Here the proof has learned that the lower prefix is not random noise.
It has an order in scale, and that order can be written as a triangular
operator.
That is the good news.
The bad news is that the triangular operator leaves affine moments and
weighted quotient modes that the available estimates do not eliminate.
Signed descent then tries to save cancellation before absolute values destroy
it.
That attempt fails at eigendirection alignment.
The quotient-kill attempts test whether the remaining affine mass can be
removed by returning to LPAS, invoking far-corona control, or spending one more
signed derivative.
Each choice returns to an earlier debt.
The proof must make that return visible.
Otherwise the reader will think the affine quotient is a new independent
obstruction instead of the lower-memory form of the same unpaid scale/source
economy.

The fourth phase is selector and projected-flow support.
Projected flow, TPS, and SG.4 put the defect inside a parabolic or two-point
equation, but that placement does not by itself supply coercivity.
The exact missing input is a selector-preserving coercive estimate, not generic
parabolic smoothing.
The paper has to preserve the exact promise and the exact failure.
The promise is not merely that a norm is controlled.
The promise is that a selected pair, direction, or defect class remains visible
long enough to produce the one-sided lower envelope needed for the route.
The failure is selector production.
Norm control can forget the very data the branch needs: sign, direction,
family membership, and one-sided growth.
That is why TPS and SG.4 belong in the failure record as downstream support
and readout pressure, not as the primitive CM proof.
They may still matter.
They can sharpen a Field readout after a same-solution packet is selected.
They cannot replace the Part and Field test unless a theorem transfers the
exact selected packet into that witness tree.

The fifth phase is the source wall.
This phase starts from the \(Q(t)\) obstruction and repeatedly tries to delete
the source.
Every serious deletion attempt has the same rhythm.
It removes legal pieces, inherited reserve, or externally accounted flux.
It then finds a smaller terminal source object that still has to be paid.
BASAC gives the clean model: \(g_m(s)=m1_{(-1/m,0]}(s)\) has
\(\|g_m\|_{L^1}=1\) and \(g_m(s)\,ds\rightharpoonup\delta_0\).
Uniform integrability, reverse Holder residence, Morrey anti-atom control, or
a Liouville theorem would pay that debt, but those are the missing supplier
theorems.
The source wall so has to be written as a productive failure.
It produces the actual objects a breakdown claim would carry.
In the CM route, those objects are admitted to the terminal test.
The proof asks whether the source object has a positive same-fluid packet,
whether the same equation participates on that packet, and whether a coherent
field readout remains at a positive scale.

The sixth phase is pressure.
Pressure has to be separated from native source because its mathematical
carrier is different.
The pressure equation is elliptic in space and tensorial in structure.
It can look source-like without having same-fluid source ancestry.
Pressure sustain so splits into source-carried pressure and pure
pressure residue.
Angular saturation sharpens the same problem directionally: global tensor
cancellation may coexist with a selected terminal lobe that keeps a positive
bill.
The paper has to make this split explicit because otherwise pressure will be
silently folded into the source wall.
That would be dishonest.
The pressure branch needs either a pressure-source tether or a pure-pressure
Liouville theorem.
Without one, it remains pressure-side support.
Its natural CM landing is Field-facing when pressure residue destroys every
coherent positive-scale field realization.  In that case full VPI
participation is lost.

The seventh phase is boundary, time, and Zeno ancestry.
Transported cylinders fix a real problem: they remove artificial boundary
flux created by the wrong frame.
But no incoming flux through the transported boundary does not rule out source
mass born inside the shrinking terminal cylinder.
Local energy has the same limitation.
It gives terminal measures and signed balances; it does not by itself forbid
the selected positive source component.
Spatial elliptic smoothing has another limitation: it can regularize space
while preserving a shrinking time slab.
The Zeno route then tries to make infinite same-fluid refill spend a quantized
resource.
Radius, heat-time depth, energy, dissipation, labels, and donor balance all
fail as uniform resources.
The resulting object is terminal zero-radius ancestry.
The CM proof has to say exactly where that object lands: usually carrier-window support
failure unless a positive-scale carrier survives, and then Part or Field must
still be tested.

The eighth phase is retained reserve and signed-pair structure.
Source reserve birth is sharper than a generic terminal source atom.
It selects the first same-fluid window where donor-square reserve is created
after inherited, legal, and residual terms have been accounted for.
That first-created reserve is difficult because it is both native and retained.
The zero-moment signed-pair branch then shows that cancellation of first moment
does not remove square reserve.
Positive and negative packets can remain on the same ledger while moments
cancel on heat subwindows.
The missing theorem is not another broad source estimate; it is a same-ledger
no-free-sink or donor-depletion statement.
The paper must write this stage with care because it is easy to overstate.
Retained reserve is not proof of exit by itself.
It is an object that must either be paid by a same-ledger theorem or routed
through the CM Part/Field test after terminal admission.

The ninth phase is the critical translator phase.
Public \(L^3\), Duhamel, terminal Leray saturation, and conditional
\(H^{1/2}\) do not all play the same role.
The local \(L^3\) route is a translator once terminal critical mass is already
on the same witness ledger.
Duhamel separates heat ancestry from nonlinear native-source response and lets
the local branch enter the CM grammar.
Terminal Leray saturation closes a signed-partner defect inside that local
family, while the independent pressure-Leray route remains open until pressure
gets its tether.
The \(H^{1/2}\) route is weaker in authority.
It sees size plus frequency placement and can suggest a pure oscillatory Field
failure, but only after same-ledger extraction and oscillation-to-Field lemmas
are installed.
The reader has to see these authority levels.
Installed translator, local closure, pressure-side open route, and conditional
critical translator are different proof objects.

The tenth phase is the final handoff.
After all positive branches have been read, the paper should leave the reader
with a small number of terminal survivor types rather than a catalogue.
The survivors are high-frequency transfer, lower-prefix active square,
weighted quotient mode, signed-pair reserve, native source atom, pure pressure
residue, Zeno ancestry, terminal Duhamel mass, and conditional critical
oscillation.
The CM routing first asks whether the selected survivor has a positive material
carrier.  That is admission support outside the CM faces.  On an admitted
record it asks whether the Navier--Stokes law still participates and whether a
coherent positive-scale field readout survives.  If participation and readout
survive, the branch is on the smooth same-solution side.
If the carrier is not admitted, the branch records carrier-window support
outside the exit conclusion.  Once admitted, direct Part failure or loss of
every positive-scale Field readout is a presentation of lost full VPI
participation.  Contraposition leaves the participating member branch smooth;
carrier support itself remains outside CM.

\subsection{Branch-To-Face Conversion Checks}

The final safeguard is to say how the branch residue would enter the class
faces without pretending that every hard object is already an exit.
The CM proof has only three primitive failures.
Every survivor below must either be support for one of those failures, a local
translator into one of them, or an open supplier theorem outside the proof.

\paragraph{High-frequency transfer.}
A retained high-frequency transfer packet first asks for Pack.
Is there still a positive transported packet carrying the same solution at the
terminal scale?
If the cascade description has no positive same-fluid carrier, the first lost
face is Pack.
If a carrier survives, the next question is Part: does the differentiated
Navier--Stokes law still act on that carrier with the same pressure and
viscosity?
Only after both questions survive does the transfer packet become Field-facing.
Then the issue is whether the active tower has one positive scale where the
field remains coherent and finite.
This prevents the scale branch from being misread as automatic smoothness or
automatic nonsmoothness.
It is a packet whose proof role depends on the first class requirement it
loses.

\paragraph{Lower-shell half-derivative deficit.}
The lower-shell deficit left by the lifted-band calculation is not itself a
class face.
It is a missing gain.
The positive route requires one of three exact inputs: an extra half derivative,
an intermediate shell gain, or a coupled weighted-tail theorem.
For Silver, this is only one possible departure presentation on the original
history.  If the deficit prevents every coherent positive-scale field
realization, the branch is Field-facing and has lost full VPI participation.
If the deficit actually destroys the carrier before participation is even
available, then the first lost face is Pack.
The paper must leave that order explicit because a half-derivative gap sounds
like a regularity failure, while the class proof needs a first failed witness
requirement.

\paragraph{Lower-prefix active square.}
The LPAS object \(P_j^\downarrow2^{-j}D_j^2\) is a support object until it is
tied to the selected terminal record.
It measures lower-prefix energy acting on the active shell.
As a positive theorem target it would become a coercive drain or a tail
payment.
As CM material it has to be located on one same-solution packet.
If the lower-prefix active square shows that no positive transported packet
survives, it supports Pack carrier-window nonadmission outside CM.
If a packet survives but the source or strain economy no longer solves the same
differentiated equation, it supports Part failure.
If the square prevents every coherent positive field scale, it is Field
support and, when realized as the terminal departure, directly exhibits
nonparticipation.
The same algebra can so have three different proof roles depending on
where the first witness requirement fails.

\paragraph{Affine quotient mode.}
The affine quotient mode is a dangerous place for overclaim.
It is small in description and persistent in effect.
The Volterra reduction shows that a large lower-prefix memory can collapse to
moments such as \(M_0\) and \(M_1\).
That does not prove those moments vanish, and it does not prove they break the
class.
The conversion check is whether the same terminal branch retains its original
carrier, law, and coherent whole-field realization.
If no carrier remains, the quotient is Pack nonadmission evidence outside CM.
If the quotient represents a same-carrier memory term that no longer
participates in the Navier--Stokes law, it is Part-facing.
If the quotient obstructs every coherent positive-scale field realization, it
is Field-facing and directly exhibits nonparticipation.
Without that first-face decision, the quotient remains a survivor, not a
conclusion.

\paragraph{Signed-pair reserve.}
Signed-pair reserve shows why cancellation has to be tied to the witness
ledger.
Zero moment can coexist with positive and negative packets whose retained
masses are bounded below by a fixed \(c_0>0\).
If the partners live on different carriers, the selected branch may fail Pack
or Part before sign even matters.
If the partners retain carrier and local equation readings, the Field question
is whether their exchange preserves the coherent whole-field realization.  If
not, full participation is lost.
The missing same-ledger no-free-sink theorem would be a positive payment.
Without it, the CM proof can still use the object only after a terminal
admission and first-Part/Field derivation.
This keeps the paper from saying that cancellation saves the branch and from
saying that cancellation failure alone proves exit.

\paragraph{Native source atom.}
A native source atom is born from the failure of temporal integrability or
anti-atom control.
The atom is terminal and source-like, but its first Part/Field is not determined by
that description alone.
If it has zero terminal radius and no positive material carrier, it is Pack
support.
If it has a carrier but the same Navier--Stokes law cannot act on that carrier
because the source participation is incompatible with pressure, viscosity, or
the selected local balance, it is Part support.
If the atom blocks every coherent positive-scale field realization, it is
Field support and, when realized terminally, direct nonparticipation.
The PDF must preserve this ordering because source atoms are the easiest place
to smuggle in a positive deletion theorem or a cheap exit label.

\paragraph{Pure pressure residue.}
Pure pressure residue is not native source with a different name.
Its conversion check starts by asking whether the pressure object is tethered
to the same-fluid packet.
Without that tether, it is support or supplier pressure, not proof material.
If the pressure residue prevents a positive transported packet from existing,
it can support Pack.
If the packet exists but pressure no longer participates correctly in the same
Navier--Stokes equation, it supports Part.
The most natural diagnostic is Field: a pure pressure lobe may destroy every
coherent whole-field realization.  If that failure is realized terminally,
full VPI participation is absent.
That is why the pressure sections repeatedly ask for pressure-source tethering
or pure-pressure Liouville.
Those are the positive doors; CM Part/Field conversion is the contrapositive door.

\paragraph{Interior time-face atom.}
No-incoming-flux arguments can remove boundary explanations for a terminal
source object.
They do not remove an atom born inside the transported cylinder.
Boundary bookkeeping only identifies this particular departure presentation;
it does not start Silver.
If the interior atom has no positive same-fluid carrier at terminal scale, the
branch is Pack-facing.
If the carrier remains but the local source balance cannot participate in the
same equation, it is Part-facing.
If the atom prevents every coherent positive field scale, it is Field-facing
and therefore nonparticipating.
This check is what keeps transported geometry honest.
The frame correction is valuable because it removes fake boundary flux.
It is not the whole theorem because it leaves an interior production question.

\paragraph{Zeno ancestry.}
Zeno ancestry is a terminal chain problem, not a fourth class face.
The conversion check is especially strict.
If the chain shrinks to zero terminal radius and no positive-scale carrier is
retained, the branch remains Pack nonadmission evidence outside CM.
If a positive-scale carrier is retained, then Zeno language alone no longer
decides the case.
The proof must ask whether the same Navier--Stokes law participates on the
retained branch.
If not, the landing is Part.
If infinite refill destroys every coherent positive-scale field realization,
the diagnostic is Field and the terminal object is not a full participant.
This is why the paper must not leave ``Zeno'' as a verdict.
It is a mechanism whose class role depends on retained carrier and
participation.

\paragraph{First-created source reserve.}
First-created source reserve is proof-relevant because it remains after past,
legal, and residual reserve have been accounted for.
The positive route wants to charge it to a same-ledger payment theorem.
The CM route asks where it sits in the witness tree.
If reserve birth happens without a positive transported packet, Pack fails.
If the packet exists but the reserve cannot be made to participate in the same
equation's pressure/viscosity/source balance, Part fails.
If reserve birth prevents a coherent readout on every positive scale, Field
fails and full participation is lost.
Until one of those happens, first-created reserve is a quarantined supplier
object.
It is too concrete to ignore and too unresolved to call a theorem payment.

\paragraph{Terminal \(L^3\) Duhamel response.}
The \(L^3\) Duhamel branch is the clean example of a translator.
Once terminal critical mass is already on the same ledger, Duhamel separates
heat ancestry from nonlinear response and routes the response into the local
CM grammar.
The branch can then support Pack, Part, legal material, or Field-facing
readout according to the local packet analysis.
Its authority is limited by that same-ledger assumption.
If the terminal mass is not on the selected witness ledger, the local
translator has not yet entered the proof.
This prevents public critical theory from being used as a generic proof of
the entire terminal branch.
It is powerful after selection; it is not a substitute for selection.

\paragraph{Critical \(H^{1/2}\) oscillation.}
The \(H^{1/2}\) branch is conditional field analysis.
Amplitude concentration can route toward the \(L^3\) translator.
Pure oscillation is different.
It may preserve amplitude while destroying frequency coherence.
If no positive coherent readout survives, that is a Field failure: the same
history may still be identifiable, but full VPI participation has been lost.
The missing pieces are same-ledger extraction and oscillation-to-Field
conversion.
Until those are proved, \(H^{1/2}\) remains a candidate translator.
The PDF must keep this conditional status because critical norms are highly
tempting places to overstate authority.

\paragraph{No fourth face.}
The conversion checks all end at the same rule.
The proof has no hidden fourth continuation requirement.
Every selected branch either lacks a positive transported packet, keeps the
packet and loses same-equation participation, keeps both and loses positive
field readout, or remains outside the class proof as support.
This is the final reason the failure record has to be so explicit.
Without the record, Part and Field can look like labels placed on
failures after the fact.
With the record, the reader sees each terminal object being produced by an
honest positive proof attempt and then tested by the three requirements that
define membership.

\subsection{Reading Rules For The Failure Record}

The reader should read this record with four distinctions kept active.
First, a route can be mathematically valuable while still failing as a direct
smoothness proof.
Second, a failed direct route can produce an object that is useful to the
contrapositive proof.
Third, an object useful to the contrapositive proof is a class exit only after
same-witness admission and a first Part/Field failure.
Fourth, no auxiliary theorem enters the proof until it has a same-witness
landing.
These distinctions are what make the back half of the paper honest.
They keep the text from treating failed estimates as embarrassing debris, and
they also keep it from treating every hard residue as if it had already proved
nonmembership.

The positive route and the CM route ask different questions.
The positive route asks for an estimate that deletes the obstruction:
control the tail, pay the source, remove the atom, tether pressure, collapse
the quotient, prevent Zeno refill, or extract a critical readout.
When such an estimate is proved, it can strengthen the pass side.
When it is not proved, the failure still teaches the paper what the terminal
object must be.
The CM route then asks a narrower question about that object.
Does it count as the same Navier--Stokes branch?
That is a membership question, not a deletion question.
The PDF must hold this distinction because otherwise every open supplier
theorem is misread as a missing proof hole in the CM argument, and every
terminal object is misread as an unlicensed shortcut.

The order of the record is also a proof device.
If the source atom appeared before the \(Q(t)\) calculation, it would look like
an invented terminal object.
If the affine quotient appeared before Volterra memory, it would look like a
name rather than a reduction.
If pressure residue appeared before source ancestry was separated from
pressure ancestry, it would look like another source term.
If \(H^{1/2}\) appeared before the installed \(L^3\) translator, it would look
like a completed critical criterion.
The ordered record prevents those mistakes.
It lets the reader watch each object become necessary.
That is why the section begins with the four-body circuit, then sharpens scale,
source, pressure, Zeno, and critical readout in sequence.

The branch table near the end compresses the preceding work; it does not
replace it.
A table can say that BASAC leaves a terminal atom, pressure leaves a pure
residue, or Volterra leaves an affine quotient.
It cannot by itself make those objects intelligible.
The preceding sections supply the reason each object exists and the exact
theorem that would have removed it.
Only after that does the table help.
The table protects the reader from losing the whole map after many branches,
but the proof value comes from the route-by-route payment analysis.
This is the difference between a written proof and a summary table.

The same rule applies to the auxiliary material behind these sections.
It can suggest definitions, route connections, and open supplier theorems.
It becomes part of the argument only after the relevant attempt, calculation,
failure, residue, and CM role have been written here in ordinary mathematical
prose.
That is the rule used in the pages that follow.

Finally, the failure record is allowed to stop short where the mathematics
really stops short.
It should not smooth over an open supplier theorem.
It should not claim a pressure tether, a temporal anti-atom theorem, a
same-ledger signed-pair theorem, a Zeno quantization theorem, or an
\(H^{1/2}\) Field extraction theorem unless that theorem is actually installed.
The honest completion of this record is not the claim that every positive
route has been closed.
It is the claim that every positive route has either paid its estimate, exposed
a terminal object, or been quarantined as support until it can land on a
selected Part or Field face.
That is exactly the kind of information the CM proof needs.

This also fixes the burden on the writer.
The writer may not use the existence of this appendix, the table of contents,
or the successful PDF build as evidence that the reader has received the
argument.
For each branch the page itself has to answer five questions without sending
the reader back to a repository surface.
What did the positive proof try to buy?
Which calculation or theorem input selected that purchase for testing?
Where did the calculation stop?
What residue survived the stopping point?
What is the residue's exact role in the Part and Field grammar?
If the page does not answer those questions, the branch has not yet been
written into the paper.
This is why the following reader-order narrative remains necessary even after
the summary packets above.
The packets give the map.
The narrative gives the mathematical experience of trying the route, watching
it stop, and carrying its survivor into the class-membership proof.
That extra experience matters because the reader has to trust the transfer of
burden.
If the paper jumps from a named failed route directly to a Part or Field
label, the reader has to supply the missing motion alone.
The text below prevents that.
It keeps the attempted proof, the failed payment, the residue, and the class
face in one line of sight, so the contrapositive proof feels earned from the
actual history of the work rather than imposed on top of it.
That continuity is part of the mathematical evidence, because it shows that the
CM test receives the same survivor the positive route actually produced.
It also gives the referee a visible audit trail through the proof transition.
That audit trail is part of the claimed proof.

\subsection{Chronological Reconstruction From Four-Body Support To CM Exit}

The history from the four-body circuit onward is not background.
It explains why the proof now has the shape it has.
The route did not begin with Part and Field as formal names.
It began with a positive regularity circuit and then learned, route by route,
that the failed positive estimates were producing terminal objects.
The class-membership proof is the point where those terminal objects stop being
treated as enemies to delete and start being tested as possible evidence of
class exit.

The chronology can be read as a sequence of changes in the live question:
\[
  \begin{aligned}
  &\text{Can the four-body circuit prove direct smoothness?}\\
  &\Longrightarrow
  \text{Which source, scale, or compactness payment is unpaid?}\\
  &\Longrightarrow
  \text{What terminal object does the unpaid payment leave?}\\
  &\Longrightarrow
  \text{Is that terminal object still the same Navier--Stokes member?}
  \end{aligned}
\]
Each arrow matters.
Removing one of them makes the final proof look like a sudden change of
strategy.
Keeping all of them shows why the final strategy is forced by the earlier
failures.

\subsubsection{Source recovery and the four-body ranking}

The first corrective event was source ranking.
The recovered material did not say that there was no proof architecture.
It said the opposite: there was too much material, at too many proof levels.
Some texts were theorem attempts.
Some were manuscript wrappers.
Some were motivation-only explanations.
Some were generated summaries of older work.
The live question so became a ranking question:
\[
  \text{which objects actually carry theorem obligations?}
\]

The answer was the four-body source stack.
The stack was not a list of four slogans.
It was a regularity circuit:
\[
  S_{\mathrm{cl}}
  \to
  Q_{\mathrm{cl}},
  \qquad
  (S_{\mathrm{cl}},Q_{\mathrm{cl}})
  \to
  C_{\mathrm{cl}},
  \qquad
  C_{\mathrm{cl}}
  \to
  G_{\mathrm{cl}},
  \qquad
  G_{\mathrm{cl}}
  \to
  S_{\mathrm{cl}}.
\]
Here \(S_{\mathrm{cl}}\) is the scale barrier, \(Q_{\mathrm{cl}}\) is the
energy/enstrophy functional, \(C_{\mathrm{cl}}\) is compact exact-object
production, and \(G_{\mathrm{cl}}\) is gradient or curvature recovery.
The order is structural.
The first arrow spends a named scale-barrier inequality
\(S_{\mathrm{cl}}\to Q_{\mathrm{cl}}\); without that inequality, scale language
does not feed the energy/enstrophy bridge.
Those controls feed compactness.
Compactness produces the object on which the gradient bridge must act.
The gradient bridge then returns control to the scale side.

The ranking step changed the proof discipline.
It demoted any surface that merely described smoothness, regularization, or
stability without paying one of these arrows.
It promoted sources that named a real payment:
high-frequency transfer, nonlinear source, compact object capture, or
same-surface derivative drain.
That is why the four-body circuit is the first real section of the failure
record.
It is the first recovered architecture whose parts can be checked as
mathematical obligations rather than as proof mood.

\subsubsection{\(Q(t)\) exposed the source wall}

The next change happened inside the second body.
The old \(Q(t)\) target was
\[
  Q(t)=\|u(t)\|_{L^2}^2+\|\nabla u(t)\|_{L^2}^2
\]
with the required conclusion that \(Q\) be a monotone or otherwise controlled
quantity.
The kinetic half did behave that way.
The enstrophy half did not.
It exposed
\[
  \int_{\mathbb T^3}(u\cdot\nabla)u\cdot\Delta u\,dx .
\]
After interpolation and the Young-margin step the proof reached
\[
  Y'(t)\le C_\nu Y(t)^3 .
\]
That inequality is not a global regularity theorem.
It is the doorway to the source wall.

This was the first decisive mutation in the proof history.
The question stopped being:
\[
  \text{can }Q(t)\text{ be called monotone?}
\]
and became:
\[
  \text{where does the nonlinear source go?}
\]
The later source-pulse, source-reserve, BASAC, pressure-sustain, terminal atom,
signed-current, and Duhamel branches all descend from that change.
They are not miscellaneous appendices.
They are attempts to answer the old \(Q(t)\) question honestly.
The source wall is the enstrophy source made terminal.

\subsubsection{The modernized packet kept the four jobs}

The modern route did not discard the four bodies.
It changed the language so the same survivor could be followed across several
coordinate systems.
The scale body became tail leakage and no-escape across dyadic bands.
The \(Q(t)\) body became source reserve and source payment.
The compactness body became exact terminal packet capture.
The curvature body became deformation geometry and commutator drain.

This modernization was useful because it prevented a false choice.
The proof did not have to decide whether the obstruction was a frequency
object, a source object, a compactness object, or a geometry object.
The correct question was whether those were several readings of one survivor.
In modern notation the intended packet has the form
\[
  \bigl(\mathcal P_{\mathrm{mod}}(T),
  \mathcal E_{\mathrm{live}}^{(n)}\bigr)
  \Longrightarrow
  \bigl(\mathcal P_{\mathrm{mod}}(T^+),
  \mathcal E_{\mathrm{live}}^{(n+1)}\bigr).
\]
The packet is successful only when the live survivor contracts, is controlled by a dissipative margin, or
is made invisible to the next scale readout.
If the survivor is merely renamed, the packet has not closed the loop.

That distinction is important for the final proof.
The manuscript may use modern language, but it cannot let modern language
pretend to be proof.
Tail leakage, source reserve, deformation commutator, and terminal packet
capture have to stay tied to the same selected branch.
The CM proof inherits this as same-record discipline:
\[
  \WitnessRecord=(u_0,\nu,u,p,Q,N,r,\Phi,T_*).
\]
No branch is allowed to switch records in the middle of the argument.

\subsubsection{Defect equations clarified the readout problem}

The projected-flow and two-point defect routes then made the survivor look like
a parabolic defect.
This was a real mathematical improvement.
A defect \(X\) with an energy identity of the form
\[
  \frac{d}{dt}E_D(X)+\mathcal D_D(X)
  =
  N_1(X)+N_2(X)
\]
would require the dissipative-margin estimate
\[
  |N_1(X)|+|N_2(X)|
  \le
  \varepsilon\mathcal D_D(X)+C_\varepsilon\Psi(E_D(X)).
\]
A two-point packet gives a similar requirement.
For a pair defect \(W_J\), the missing statement is a coercive drift-diffusion
estimate whose controlled remainder implies coherence of the pair readings.

The failure was not that defect equations were useless.
The failure was selector readout.
The needed one-sided statement has the shape
\[
  \frac{d}{dt}\log d_{ab}^{\mathrm{coarse}}(t)
  =
  s_J(a,b,t),
  \qquad
  s_J(a,b,t)\ge \lambda_J(t)-\varepsilon_J(a,b,t).
\]
The exact missing input is a one-sided directional lower-envelope theorem for
the selected pairs; a norm estimate on a defect does not supply it.
It may control size while losing sign, direction, selected-pair membership, and
one-sided expansion.

The branch splits there.
Defect equations can guide the proof.
They do not become the proof primitive.
The primitive had to be the selected same-solution terminal object and its
readout.
This is one reason the final grammar contains \(\Field_{N,r,Q}\).
Field records whether the same history retains coherent positive-scale field
realization; its failure is loss of full participation.

\subsubsection{Source-wall deletion produced sharper survivors}

The source-wall phase repeatedly tried to delete the bad source object.
The pattern was stable:
\[
  \begin{aligned}
  &\text{identify a source pulse, reserve, donor object, or current}\\
  &\Longrightarrow
  \text{try to charge it to a native ledger}\\
  &\Longrightarrow
  \text{find a sharper survivor}\\
  &\Longrightarrow
  \text{try to charge that survivor.}
  \end{aligned}
\]
This was not wasted work.
Each cycle made the terminal object more precise.

The model time atom explains the wall.
For
\[
  g_m(s)=m\,\mathbf 1_{(-1/m,0]}(s),
\]
one has
\[
  \int_{-1}^{0}g_m(s)\,ds=1,
  \qquad
  \int_{-1}^{0}g_m(s)^p\,ds=m^{p-1}\to\infty
  \quad(p>1).
\]
Thus the fixed \(L^1_t\) bound and weak convergence to \(\delta_0\) do not give
uniform integrability.
The deletion route then needs an anti-atom theorem, a reverse Holder residence
theorem, a Morrey decay theorem, a rigidity theorem, or a Liouville theorem.
Without one of those, the terminal source atom remains.

Pressure did the same thing in a different language.
Pressure can split into a source-carried part and a singular pressure part
orthogonal to native source ancestry:
\[
  \mu_*^{\mathrm{press}}
  =
  f\,\Lambda_*^{\mathrm{src}}
  +
  \mu_*^{\mathrm{press}\perp\mathrm{src}}.
\]
The first term can join the source account.
The second term demands a pressure-source tether or a pure-pressure Liouville
theorem.
Again the proof did not obtain deletion for free.
It obtained a sharper survivor.

This is the lived source-wall lesson.
Positive proof work can make an obstruction more exact without erasing it.
The final proof must not throw that exactness away.
It has to ask whether the exact survivor is still a lawful terminal
Navier--Stokes member.

\subsubsection{The bad object changed proof role}

The deepest change was the role of the bad object itself.
In the positive route, the bad object was the enemy:
\[
  \text{bad object exists}
  \Longrightarrow
  \text{prove it cannot exist}
  \Longrightarrow
  \text{smoothness}.
\]
That route is legitimate when the deletion theorem is available.
It becomes theorem-sized when every branch asks for a new deletion theorem.

The class-membership route changes the role:
\[
\begin{aligned}
&\text{bad terminal object exists}\\
&\Longrightarrow\ \text{test it as the same terminal object}\\
&\Longrightarrow\ \text{derive the first Part/Field failure}\\
&\Longrightarrow\ \Exit(Q;\Owork).
\end{aligned}
\]
The bad object becomes an exit witness only after it enters the terminal CM
test.
But once it is admitted and a first Part/Field failure is derived, it is no longer an
enemy of the proof.
It is the witness that the alleged terminal object is outside the class.

This corrected two opposite errors.
The first error was the cheap shortcut:
\[
  \text{bad object}\Longrightarrow\Exit(Q).
\]
That is false without terminal admission and face typing.
The second error was the positive overburden:
\[
  \text{every bad object must be deleted before the proof can close.}
\]
That is also false for a contrapositive class-exit proof.
Some bad objects are exactly what the fail branch is supposed to produce.

\subsubsection{Terminal Clay entry fixed the arrow}

The terminal Clay branch then fixed the direction of the argument.
The paper cannot begin by saying that class exit makes a Clay witness
inadmissible.
That skips the counterexample's own claim.
The honest order starts with the alleged finite terminal witness:
\[
  \begin{aligned}
  &\text{finite Clay terminal witness}\\
  &\Longrightarrow
  \text{same original datum, equation, fluid, terminal role}\\
  &\Longrightarrow
  \text{same-fluid CM terminal participation-field tree}\\
  &\Longrightarrow
  \text{first failed Part/Field face}\\
  &\Longrightarrow
  \Exit(Q;\Owork).
  \end{aligned}
\]
This is the entry-first correction.
It treats the alleged breakdown object seriously enough to test it.
It also prevents the proof from rejecting the object by definition.

The Field-certification questions are now unavoidable:
\[
  \begin{array}{ll}
  \Pack_Q: & \text{does a positive same-fluid carrier survive?}\\
  \Part_{N,Q}: & \text{does the same pressure-viscosity tower act there?}\\
  \Field_{N,r,Q}: & \text{does one positive scale read one coherent field?}
  \end{array}
\]
These are not labels pasted onto the failure record.
They are the terminal entry test demanded by the history of the failed
positive routes.

\subsubsection{The pass/exit engine}

The final proof behavior is the pass/exit engine.
For each obstruction \(O\), the branch split is
\[
  O_{\Pass}
  \qquad\text{or}\qquad
  O_{\Fail}.
\]
The pass branch keeps the same-fluid requirement:
\[
  O_{\Pass}
  \Longrightarrow
  \Part_{N,Q}\wedge\Field_{N,r,Q}
  \Longrightarrow
  \Member(Q;\Owork)
  \Longrightarrow
  \text{smooth continuation}.
\]
The fail branch is not a loose bad scenario.
It either remains at the carrier-window support boundary before admission, or it
is a terminal object capable of finite nonsmoothness and admitted to the CM test:
\[
  O_{\Fail}
  \Longrightarrow
  \neg\Part_{N,Q}
  \vee
  \forall r>0\,\neg\Field_{N,r,Q}.
\]
Then and only then does the proof read
\[
  \Exit(Q;\Owork):=\neg\Member(Q;\Owork).
\]

The no-third-branch law is what makes this close.
There is no legal branch that is both an in-class member continuation and the
same terminal Part/Field-failure branch.
So the obstruction is either the smooth member side or the class-exit side.
It is not a third in-class nonsmooth continuation.

\subsubsection{The source-reserve stress test}

The source-reserve burst is the best stress test of this logic because it
produced a real surviving branch.
Choose the first terminal same-fluid window \(W\) where donor-square reserve is
born, and write
\[
  R_N(W)
  =
  \int_W
  \sum_{k>N}2^k
  \left(\sum_{\ell>k+4}D_\ell(t)\right)^2\,dt.
\]
Minimality removes inherited past reserve.
Legal and residual ledgers remove legal and child terms.
The survivor is first-created positive-scale native donor-square reserve.

The zero-moment signed-pair branch then shows why visibility is hard.
One can have
\[
  J=J_+-J_-,
  \qquad
  \int_W J_+ = \int_W J_-,
\]
while both \(J_+\) and \(J_-\) have retained mass at least \(c_0>0\) on the
same retained terminal ledger.
The first moment cancels.
The reserve does not disappear.
This is a genuine forward-positive blocker for a no-free-sink theorem.

The CM reading is more precise.
The branch has four possible roles.
It may be a pass-side requirement.
It may be admitted to the terminal CM test.
It may produce a Part or Field Part/Field failure.
Or it may remain supplier quarantine until the same-ledger payment theorem is
proved.
The existence of the signed-pair gap does not by itself reopen the CM route.
To reopen the route, it would have to produce a genuine third branch:
a finite terminal object that remains in \(\Member(Q;\Owork)\) and also blocks
smooth continuation, or a failure of terminal CM entry, finite failure
exhaustion, or contrapositive embedding.

\subsubsection{Why one frontier became target topology}

The final chronology lesson is that there is no single live frontier that can
stand for the whole proof.
At least five proof modes have to stay separated:
\[
  \begin{gathered}
  \text{positive smoothness estimates},\\
  \text{CM contrapositive class exit},\\
  \text{source-wall supplier work},\\
  \text{receiver, readout, and endpoint support},\\
  \text{comparison or transfer work.}
  \end{gathered}
\]
A theorem can be decisive in one mode and only supportive in another.
A source-reserve visibility theorem is a forward-positive target.
It is CM closure only after it lands on the same selected terminal witness and
derives the relevant Part or Field consequence.
An endpoint readout theorem is pass-side support after Part and Field
survive.
It is not the upstream class test.
Euler is admitted only by a forward limit from the original Navier--Stokes
history in which effective viscosity and every forcing, Reynolds-stress,
pressure, incompressibility, compactness, and ancestry defect vanish.  If any
residual remains, the destination is the corresponding residual law rather
than Euler.  In every case the comparison only describes loss of the original
fixed-\(\nu\) participation.

The manuscript so treats topology as part of the proof.
Each object has to be placed by proof role before it can be used.
The current reader-facing rule is:
\[
  \text{same object, same record, same proof mode, same face order.}
\]
This rule is what prevents the PDF from turning the recovered source field into
a pile of named surfaces.
It also prevents the opposite failure: compressing the whole route into one
frontier label that hides the different burdens.

\subsubsection{Chronology-to-proof handoff}

The chronology from the four-body circuit to the CM engine can now be stated as
a proof handoff.
The four-body route supplies the original positive payment scheme.
The \(Q(t)\) correction exposes the source wall.
The modern packet identifies one survivor through tail, source, compactness,
and deformation readings.
The defect routes show that norm control is weaker than the needed one-sided
readout.
The source-wall routes sharpen source, pressure, atom, Zeno, and reserve
survivors.
The CM correction changes those survivors from objects to delete into objects
to test.
The terminal Clay entry rule makes the alleged finite breakdown witness enter
the same-fluid participation-field tree.
The pass/exit engine then consumes the branch.

This handoff is what the rest of the appendix is doing in detail.
Every later branch has the same reader-facing form:
\[
  \begin{aligned}
  &\text{positive attempted payment}\\
  &\Longrightarrow
  \text{calculation or theorem input that selected the payment for testing}\\
  &\Longrightarrow
  \text{exact stopping point}\\
  &\Longrightarrow
  \text{surviving terminal object}\\
  &\Longrightarrow
  \text{Part and Field, pass-side readout, or supplier quarantine.}
  \end{aligned}
\]
That is not a status format.
It is the mathematical bridge from the source chronology to the final
contrapositive proof.

\subsection{The Attempt Chain In Reader Order}

The record begins with the four-body circuit because that is where the
smoothness project first stopped being a general hope and became a payment
scheme.
The circuit did not merely say that viscosity should smooth the fluid.
It tried to make four statements pay one another.
The route required four exact statements.
The scale statement controls high-frequency energy.
The \(Q(t)\) statement binds energy and enstrophy into one controlled quantity.
The compactness statement selects an exact Navier--Stokes object from controlled
approximations.
The gradient statement recovers the derivative control needed to return to the
scale statement.
The proposed proof was circular in the intended good sense: the four statements
form a closed continuation loop once each statement is proved on the same
periodic surface.

That shape explains both the route and the failure.
Each body names a familiar mathematical move, but each move has an exact input.
High-frequency dissipation is real.
The kinetic energy identity is exact.
Compactness can pass from approximations to a solution only after the topology,
time-compactness, uniform tail smallness, nonlinear passage, and same-solution
attachment have all been supplied.
Derivative recovery is exactly what a smoothness proof must eventually provide.
The failure is that none of those bodies can be allowed to borrow the theorem
from the next body.
The scale body needs a nonlinear transfer estimate before it can suppress the
tail.
The \(Q(t)\) body needs the vortex-stretching source to be paid before it can be
monotone in the required sense.
The compactness body needs the first two controls before it can produce the
right limit.
The gradient body, as recovered from the source stack, used geometric curvature
language; the missing input is a periodic flat-torus derivative-recovery theorem
on the same surface.
The circuit so identifies the correct proof architecture and also shows
where the old proof overclaimed.

The first body, the scale barrier, tried to say that the cascade cannot keep
moving energy into arbitrarily high modes faster than viscosity removes it.
In spectral language it starts from
\[
  \partial_t E(k,t)=T(k,t)-2\nu\,k^2E(k,t)+F(k,t),
\]
and on the zero-force problem \(F=0\).
The required positive-route estimate is a transfer-tail bound such as
\[
  T(k,t)\le Ck^{-\alpha}
\]
with enough decay to leave the viscous term dominant at high \(k\).
Combined with a spectral decay law \(E(k,t)\le Ck^{-\beta}\), \(\beta>3\), this
gives
\[
  \|\nabla u(t)\|_{L^2}^2
  \le
  C_{\mathrm{spec}}\int_0^\infty k^2E(k,t)\,dk
  <\infty .
\]
This is why the route looked like a proof.
It replaces the old sentence ``viscosity wins at small scales'' by an actual
summability condition.
The exact point of failure is the transfer law.
The energy identity pays \(\int\|\nabla u\|_2^2\) in time.
It does not supply the pointwise cascade decay needed to force
\(T(k,t)\le Ck^{-\alpha}\) with the exponent that makes the derivative tail
summable.
The residue left by this failed attempt is not vague.
It is the high-frequency nonlinear transfer packet.
Later names such as lifted-band leakage, gap kernels, tail stress, and active
square all descend from this first unpaid scale bill.

The second body, \(Q(t)\), tried to make the energy law carry one more
derivative.
The functional was
\[
  Q(t)=\|u(t)\|_{L^2}^2+\|\nabla u(t)\|_{L^2}^2 .
\]
The kinetic half works:
\[
  \frac12\frac{d}{dt}\|u(t)\|_{L^2}^2
  +\nu\|\nabla u(t)\|_{L^2}^2=0 .
\]
The enstrophy half exposes the actual obstruction.
Writing \(Y(t)=\frac12\|\nabla u(t)\|_{L^2}^2\) and
\(Z(t)=\|\Delta u(t)\|_{L^2}^2\), one gets
\[
  \frac{dY}{dt}+\nu Z
  =
  \int_{\mathbb T^3}(u\cdot\nabla)u\cdot\Delta u\,dx .
\]
Holder, Sobolev interpolation, and a Young-margin split give
\[
  \left|\int (u\cdot\nabla)u\cdot\Delta u\,dx\right|
  \le
  C_S C_I
  \|\nabla u\|_{L^2}^{3/2}\|\Delta u\|_{L^2}^{3/2},
\]
and after the fixed dissipative fraction is spent it leaves
\[
  \frac{dY}{dt}\le C_\nu Y^3 .
\]
That inequality is enough for local control.
It is also compatible with finite-time growth.
So the \(Q(t)\) route fails precisely where the nonlinear source has to be
paid.
The failed route teaches the later proof what to look for: a terminal source,
reserve, or pressure object is not a side issue; it is the old \(Q(t)\)
obstruction made explicit.

The third body, non-Sobolev compactness, tried to pass from controlled
approximations to the exact object.
It was built from a tailored norm combining low-order energy with a spectral
tail:
\[
  \|u\|_{H_{\mathrm{ns}}}
  =
  \|u\|_{L^2}
  +\|\nabla u\|_{L^2}
  +\sup_{k>k_0}\frac{E(k)}{k^\beta}.
\]
The mechanism is honest.
Split \(u_n\) into low and high frequencies.
The low-frequency part is finite-dimensional once time compactness is supplied.
The high-frequency part is small when the spectral tail is small.
Strong \(L^2\) convergence then passes the quadratic term
\[
  u_n\otimes u_n\to u\otimes u .
\]
The failure is that compactness is downstream.
It does not create the scale tail or the enstrophy control.
It consumes them.
The compactness attempt so survives as a precise packet-capture lesson:
once a terminal or prelimit object is controlled in the right way, compactness
can make it exact; compactness cannot be used to pretend the missing controls
have already been proved.

The fourth body, curvature or gradient recovery, tried to close the loop by
recovering derivative control from geometric structure.
The recovered source equation is not a flat-torus gradient estimate.
It is a Riemannian reformulation,
\[
  \partial_t u+\nabla_u u
  =
  -\operatorname{grad}_g p+\nu\Delta_g u,
\]
with the model gradient inequality
\[
  \partial_t|\nabla u|^2
  \le
  -\operatorname{Ric}(g)(\nabla u,\nabla u)
  +\nu\Delta_g|\nabla u|^2
  +C_g|\nabla u|^3 .
\]
The reason this looked powerful is clear.
A positive Ricci term would act like a derivative drain.
That is exactly the fourth payment the circuit needed.
The failure is also clear.
The periodic Clay problem lives on the flat torus.
The Ricci drain does not appear there as a free same-surface theorem.
The residue is the correct demand for a same-surface replacement: either a
periodic gradient-transfer theorem, or a deformation-geometry theorem that
derives the needed drain from the actual transported Navier--Stokes structure.

The modernized four-body packet kept the four jobs and stopped pretending that
the old names were enough.
The scale tail became tail leakage.
The \(Q(t)\) source became source reserve.
The compactness object became terminal packet capture.
The geometric drain became deformation geometry and commutator control.
This was progress because it put the same survivor under several lenses instead
of treating every new name as a new proof.
The packet became a way to ask whether the dyadic tail, the derivative tower,
and the moving-coordinate geometry were the same unpaid object.
The failure remained a contraction failure.
The modern packet redescribed the survivor more accurately, but it did not yet
prove a same-surface dissipative-margin law strong enough to send Body IV back to Body I.

The kernelized lifted-band route is the first fully scale-local descendant of
the scale barrier.
It attacks the off-diagonal shell region \(N+M<k<j-4\) and reaches the carrier
\[
  \sum_{j\ge N}\sum_{N+M<k<j-4}
  2^{-(j-k)}2^{3k/2}\|\Delta_k u\|_2
  \|\nabla\Delta_j u\|_2\|\Delta_j u\|_2 .
\]
The factor \(2^{-(j-k)}\) is a real gain.
It says that separated shells do not interact with full strength.
After the Young-margin split, the route leaves a lower-shell remainder with an extra
factor \(2^k\).
That factor is the point where the estimate stops.
The installed energy surface controls \(2^{2k}\|\Delta_k u\|_2^2\) in time, not
\(2^{3k}\|\Delta_k u\|_2^2\).
The residue is so a half-derivative or equivalent lower-shell gain
target.
In CM language this branch does not become an exit by being hard.
It becomes useful after a terminal witness is selected and the retained shell
record is tested for positive-scale Field coherence.

The shell-gain branch tried to pay that extra lower-shell factor directly.
It introduced
\[
  A_j(t)=\sum_{k\le j}2^{5k/2}\|\Delta_k\omega(t)\|_2,
  \qquad
  B_j(t)=\frac{A_j(t)}{\nu 2^{2j}} .
\]
The intended conclusion was that \(B_j(t)\) is small on active windows.
The exact needed consequence is control of low-frequency strain as a coefficient
bounded by the viscous scale \(2^{2j}\).
The exact scale-matching point is:
the high shell dissipates on time \(2^{-2j}\), so lower-shell strain has to be
measured against \(\nu 2^{2j}\).
The failure is that the available dissipation is an integrated square budget.
It does not give the active-window pointwise suppression
\[
  A_j(t)\le \varepsilon\nu 2^{2j}.
\]
The residue is a dynamic suppression theorem for low-frequency strain.
Until that theorem is proved, the half-derivative target remains a forward
positive target, not a completed smoothness proof.

Cumulative tail stress and LPAS tried to move from isolated shell algebra to a
whole tail payment.
The continuum object sees a lower scale \(\ell\), a larger scale \(r\), and a
kernel factor \(\ell/r\).
Its dyadic shadow is the lower-prefix active square
\[
  \mathrm{LPAS}_N
  =
  \int_0^T\sum_{j\ge N}
  \mathcal P_j^\downarrow(t)\,2^{-j}D_j(t)^2\,dt,
\]
where \(\mathcal P_j^\downarrow\) is the energy below the active shell.
The exact reduction packages the whole lower prefix instead of chasing one shell
at a time.
It fails because the orientation is wrong for the estimates already installed.
Upper-tail queues control shells above the active index.
LPAS needs lower-prefix energy multiplied by the square of the active shell.
The crude bound by \(\|\nabla u\|_2^2\) leaves the active square uncontrolled.
The residue is exactly
\[
  \mathcal P_j^\downarrow\,2^{-j}D_j^2 .
\]
This object later reappears as active-square, source-reserve, and lower-memory
pressure.

The Volterra route tried to use the triangular order in scale.
In logarithmic scale \(s\), the lower prefix becomes a Volterra operator
\[
  (\mathcal V_\kappa f)(s)=\int_{s_0}^{s-\kappa}f(\sigma)\,d\sigma .
\]
That move matters because a lower prefix is not an arbitrary bilinear error.
It has memory.
The route splits the output into an exact-in-scale part and a persistent
weighted quotient mode.
The quotient is represented by the affine moments
\[
  M_0(t)=\int_I f(s,t)\,ds,
  \qquad
  M_1(t)=\int_I s\,f(s,t)\,ds .
\]
This was one of the cleanest reductions.
The large lower-prefix object became a small number of scale moments.
The failure is that the weight can degenerate and the inherited estimates do not
force those moments to vanish.
The residue is the affine quotient mode.
For the later class-membership proof, that residue becomes a Part/Field
Part/Field failure only after it is attached to the selected terminal object and tested
for Part and Field.

Signed descent tried to kill the Volterra quotient without throwing away
cancellation.
The key idea was to keep the scale derivative
\[
  Q_\ell=-\ell\partial_\ell P_{\le \ell}
\]
visible.
The attempted signed-commutator identity target was
\[
  \mathfrak c_\ell
  =
  (\ell\partial_\ell)\Psi_\ell
  +\nabla_x\cdot J_\ell
  +\partial_t R_\ell
  +\mathcal S_\ell^{\mathrm{spill}}
  +\mathcal E_\ell .
\]
If the scale derivative telescopes and the spill is controlled, the bad
positive majorant should never be needed.
The failure is the sign.
The principal lifted commutator contains a strain form with no fixed sign.
Active wave packets can align with expanding eigendirections of the low shell.
The route so stops at a real missing theorem: eigendirection
decorrelation, signed measure cancellation, or a nonpositive divergence
identity.
Taking absolute values sends the proof back to the active-square wall.

The projected-flow, TPS, and SG.4 routes tried to put the defect into a
parabolic or sidecar equation where diffusion could be made coercive.
The projected-flow sequence was \(D.1\) through \(D.5\): well-posedness,
coercive energy, exact lift, continuation, and final Clay integration.
The energy identity had the desired form
\[
  \frac{d}{dt}E_D(X)+\mathrm{dissipation}
  =
  N_1(X)+N_2(X),
\]
with the target estimate
\[
  |N_1|+|N_2|
  \le
  \varepsilon\,\mathrm{dissipation}+C_\varepsilon\Psi(E_D).
\]
The two-point strain route has the parallel exact target: make a defect equation
for pair strain and prove that selected pairs renew the lower envelope needed by
SG.4.
The failure is selector production.
A norm can bound the size of the defect and still lose sign, direction, pair
membership, and one-sided growth.
Thus the projected-flow package and TPS estimates are support until they prove
the exact one-sided selector theorem or land a selected same-solution packet in
the CM witness tree.

The source-wall deletion phase began when the old \(Q(t)\) failure was read as
a terminal source problem.
The repeated pattern was:
find a source pulse, charge it to a native ledger, discover a sharper survivor.
The BASAC branch is the clearest model.
The time density
\[
  g_m(s)=m\,\mathbf 1_{(-1/m,0]}(s)
\]
has \(\|g_m\|_{L^1}=1\) and
\(g_m(s)\,ds\rightharpoonup\delta_0\), while
\[
  \int_{-1}^0g_m(s)^p\,ds=m^{p-1}\to\infty
\]
for every \(p>1\).
That one calculation explains the whole wall.
An \(L^1_t\) source account can carry terminal mass without giving uniform
integrability.
The missing theorem is an anti-atom theorem: temporal uniform integrability,
reverse Holder residence, Morrey decay, rigidity, or a Liouville mechanism.
The later CM route uses this differently.
It does not need to delete the terminal atom by forward estimate.
It tests the atom as a terminal obstruction and asks which witness it loses
first.

The pressure branches show why not every source-like object is the same source.
Pressure sustain decomposes into a part absolutely continuous with native source
and a singular pressure part
\[
  \mu_*^{\mathrm{press}\perp\mathrm{src}} .
\]
Calderon--Zygmund structure relates pressure to quadratic gradients on each
time slice.
The exact missing input is transported same-fluid pressure-source ancestry for
the selected pressure contribution.
Angular saturation has the same problem in directional form.
The trace of the pressure Hessian may balance over directions, while the
terminal ledger keeps one selected positive lobe with its own weights and
filters.
So the pressure routes leave a fork: prove a pressure-source tether, or prove a
pure pressure Liouville theorem.
Until then, pressure is a pressure-side diagnostic and possible Field support,
not native source deletion.

Transported boundary and Zeno ancestry tried to remove source residues by
following the same fluid.
Transported cylinders correctly eliminate fake fixed-frame boundary inflow.
They do not eliminate an interior time-face atom such as
\[
  F_N^{\mathrm{src},+}(t)
  =
  \tau_N^{-1}\mathbf 1_{[T-\tau_N,T]}(t).
\]
The mass can be born inside the transported cylinder as the terminal layer
shrinks.
The Zeno ancestry route then tries to make infinitely many same-fluid refills
spend a quantized resource.
The failure is that all natural resources can split down the terminal scale
tree: radii shrink, heat-times sum, energy divides, and donor balances telescope
only over finite truncations.
What survives is the terminal zero-radius ancestry object.
In the later witness tree, zero-radius ancestry is carrier-window support unless a retained
positive-scale carrier is separately produced.

The May source-reserve burst sharpened the source wall into a retained reserve
problem.
The first terminal same-fluid window \(W\) carries a reserve
\[
  R_N(W)
  =
  \int_W\sum_{k>N}2^k
  \left(\sum_{\ell>k+4}D_\ell(t)\right)^2\,dt .
\]
Minimality removes inherited past reserve and leaves the first-created native
reserve.
That is the real survivor:
\[
  \text{first-created positive-scale native donor-square reserve}.
\]
The zero-moment signed-pair attempt then shows why visibility is hard.
One can have square reserve bounded below by a fixed \(c_0>0\) and zero first
moment:
\[
  J=J_+-J_-,
  \qquad
  \int_WJ_+=\int_WJ_- .
\]
The positive and negative packets can survive on the same retained ledger while
their first moments cancel on heat subwindows.
The missing theorem is a same-ledger signed-pair no-free-sink theorem.
In CM terms, the retained reserve remains supplier-quarantined until a
same-ledger payment theorem is proved or the terminal object is admitted to the
CM test and a first Part/Field failure is derived.

The critical \(L^3\), Duhamel, Leray, and \(H^{1/2}\) branches are readout
branches rather than fresh positive smoothness proofs.
The \(L^3\) route starts with terminal critical mass and uses Duhamel:
\[
  u_m(t)
  =
  e^{(t-t_0)\nu\Delta}u_m(t_0)
  +
  \int_{t_0}^{t}
  e^{(t-s)\nu\Delta}
  P\nabla\cdot(u_m\otimes u_m)(s)\,ds .
\]
The heat side is ancestry.
The nonlinear side is native source response.
Once the branch is on the same ledger, the local translator routes it into
Pack, Part, legal, or Field-facing readout.
Terminal Leray saturation handles the local signed-partner defect in the same
family.
The independent pressure-Leray route remains open because pressure still needs
its tether.
The \(H^{1/2}\) branch is conditional: amplitude routes toward \(L^3\), while
pure oscillation remains Field-facing until same-ledger extraction and
oscillation-to-Field lemmas are proved.

This is the reason the failure record belongs inside the reader-facing proof,
alongside the CM payoff.
The positive attempts are not embarrassing debris.
They identify the exact terminal objects a finite-time breakdown claim would
have to carry.
The later proof does not erase those objects.
It asks whether the same object still has a positive packet, whether the same
Navier--Stokes law still acts on that packet, and whether the resulting field
has a coherent positive-scale readout.
Those are the Part and Field questions.
The pass branch is the branch that keeps all three requirements and so
stays on the smooth same-solution side.
The fail branch is useful only after it is admitted to the terminal test and
loses one of those requirements.
That is why a source atom, a pressure residue, a Zeno chain, an affine quotient,
or a critical oscillatory packet is not called class exit merely because it is
bad.
It becomes class exit only when the first failed witness has actually been
derived.

\subsection{The Four-Body Attempt}

The first source-grounded program was a four-body circuit.
It tried to move one packet of control around a loop:
\[
  S_{\mathrm{cl}}\longrightarrow Q_{\mathrm{cl}},
  \qquad
  (S_{\mathrm{cl}},Q_{\mathrm{cl}})\longrightarrow C_{\mathrm{cl}},
  \qquad
  C_{\mathrm{cl}}\longrightarrow G_{\mathrm{cl}},
  \qquad
  G_{\mathrm{cl}}\longrightarrow S_{\mathrm{cl}} .
\]
The four bodies were:
\[
  S_{\mathrm{cl}}
  =
  \text{scale barrier},
  \qquad
  Q_{\mathrm{cl}}
  =
  \text{energy/enstrophy functional},
\]
\[
  C_{\mathrm{cl}}
  =
  \text{compactness / exact-object production},
  \qquad
  G_{\mathrm{cl}}
  =
  \text{gradient or curvature recovery}.
\]

The route looked natural.
Trapped high-frequency energy leaves the derivative stack with no place to
hide a blow-up.
If energy and enstrophy stay under a useful bound, the nonlinear transfer budget
is constrained.
If those controls are compact enough, an approximation sequence has to converge
to an exact Navier--Stokes object.
If the exact object has derivative control, the scale barrier can be certified
again.

The loop would have proved a propagation statement:
\[
  \mathcal P_{\mathrm{cl}}(T)
  =
  (S_{\mathrm{cl}}(T),Q_{\mathrm{cl}}(T),C_{\mathrm{cl}}(T),G_{\mathrm{cl}}(T))
  \quad\Longrightarrow\quad
  \mathcal P_{\mathrm{cl}}(T^+)
\]
for a later time \(T^+>T\).
That is the real shape of the first program.
The proof had to pay each arrow strongly enough to keep the same packet moving
on the same solution class.

The circuit failed because each arrow secretly needed the next hard theorem.
The scale barrier needed a nonlinear transfer bound.
The \(Q(t)\) body needed the vortex-stretching source to be paid.
Compactness needed the scale and energy inputs before the limit was useful for
the same-packet argument.
The fourth body in the recovered source stack used geometric curvature language,
while the Clay problem on \(\mathbb T^3\) needs a same-surface periodic theorem.

\subsection{Body I: Scale Barrier}

The scale-barrier attempt starts from the spectral balance
\[
  \partial_t E(k,t)=T(k,t)-2\nu\,k^2E(k,t)+F(k,t).
\]
On the zero-force branch \(F=0\).
The viscous term has the right sign at high frequency.
If the nonlinear transfer satisfies
\[
  T(k,t)\le Ck^{-\alpha},\qquad \alpha>2,
\]
and if the spectrum satisfies
\[
  E(k,t)\le Ck^{-\beta},\qquad \beta>3,
\]
then the first derivative stays finite:
\[
  \|\nabla u(t)\|_{L^2}^2
  \le
  C_{\mathrm{spec}}\int_0^\infty k^2E(k,t)\,dk
  \le
  C_{\mathrm{spec}}C\int_0^\infty k^{2-\beta}\,dk
  <\infty .
\]

The exact proof burden is not the visual appearance of a critical scale.  For
\(K>0\), define the high-frequency enstrophy tail and the
viscous-minus-positive-transfer dominance functional
\[
  Y_{>K}(t):=\int_K^\infty k^2E(k,t)\,dk,
  \qquad
  \mathcal D_K(t):=
  2\nu\int_K^\infty k^4E(k,t)\,dk
  -
  \int_K^\infty k^2T_+(k,t)\,dk,
\]
where \(T_+=\max(T,0)\).  The spectral balance gives the exact implication
\[
  \frac{d}{dt}Y_{>K}(t)
  \le
  -\mathcal D_K(t).
\]
A drain theorem therefore has to prove \(\mathcal D_K(t)\ge0\), or an
integrable bound on its negative part, on the same high-frequency tail being
used in the compactness or continuation argument.  The power-law calculation
\(T(k,t)=C_Tk^{-\alpha}\), \(E(k,t)=C_Ek^{-\beta}\) is only a profile-restricted
calculation: under those extra profile assumptions, the balance equation gives
\[
  k_{\mathrm{crit}}
  =
  \left(\frac{C_T}{2\nu C_E}\right)^{1/(2+\beta-\alpha)}.
\]
That model critical scale is not a proof for an arbitrary terminal spectrum.

The failure is the transfer estimate.
The energy identity controls
\[
  \int_0^T\|\nabla u(t)\|_{L^2}^2\,dt,
\]
and the total kinetic energy.
The high-frequency dominance theorem remains open:
\[
  \mathcal D_K(t)\ge0
  \quad\text{or}\quad
  \int_0^{T_*}(\mathcal D_K(t))_-\,dt<\infty
\]
with constants and cutoff \(K\) tied to the same selected tail.
That dominance law is already a cascade-control theorem.
Once the proof asks for it, the proof has moved the hard part into the tail.

What survived is the right object: high-frequency nonlinear transfer.
Later work reappears under sharper names:
lifted-band leakage, lower-prefix active square, Volterra scale memory,
kernelized tail stress, and source reserve.

\subsection{Body II: The \(Q(t)\) Functional}

The second body tried to bind kinetic energy and enstrophy in one named
functional:
\[
  Q(t)
  =
  \|u(t)\|_{L^2}^2+\|\nabla u(t)\|_{L^2}^2 .
\]
The kinetic-energy half is exact:
\[
  \frac12\frac{d}{dt}\|u(t)\|_{L^2}^2
  +\nu\|\nabla u(t)\|_{L^2}^2
  =0 .
\]
The enstrophy half is where the problem enters.
Let
\[
  Y(t)=\frac12\|\nabla u(t)\|_{L^2}^2,
  \qquad
  Z(t)=\|\Delta u(t)\|_{L^2}^2 .
\]
After differentiating the equation and integrating by parts,
\[
  \frac{dY}{dt}+\nu Z
  =
  \int_{\mathbb T^3}(u\cdot\nabla)u\cdot\Delta u\,dx .
\]
H\"older's inequality gives
\[
  \left|
    \int_{\mathbb T^3}(u\cdot\nabla)u\cdot\Delta u\,dx
  \right|
  \le
  \|u\|_{L^6}\|\nabla u\|_{L^3}\|\Delta u\|_{L^2}.
\]
Let \(C_S\) be the Sobolev constant and \(C_I\) the interpolation constant on
\(\mathbb T^3\).  By Sobolev and interpolation,
\[
  \|u\|_{L^6}\le C_S\|\nabla u\|_{L^2},
  \qquad
  \|\nabla u\|_{L^3}
  \le C_I
  \|\nabla u\|_{L^2}^{1/2}\|\Delta u\|_{L^2}^{1/2}.
\]
Hence
\[
  \left|
    \int_{\mathbb T^3}(u\cdot\nabla)u\cdot\Delta u\,dx
  \right|
  \le C_SC_I
  \|\nabla u\|_{L^2}^{3/2}\|\Delta u\|_{L^2}^{3/2}.
\]
Young's inequality spends a fixed part of \(\|\Delta u\|_{L^2}^2\) and leaves
\[
  \frac{dY}{dt}
  \le
  C_\nu Y^3 .
\]

This is the first familiar wall.
The inequality is useful for local theory.
It also permits the ODE blow-up profile
\[
  Y'(t)=C_\nu Y(t)^3,
  \qquad
  Y(t)
  =
  \frac{Y(0)}{\sqrt{1-2C_\nu Y(0)^2t}} .
\]
The calculation so shows exactly what the \(Q(t)\) route would have to
improve.
It needs the nonlinear source to be paid by something stronger than the
kinetic-energy identity and the enstrophy differential inequality displayed
above.

The source-forensic correction says the same thing in the old notation.
The clean dissipative \(Q(t)\) formula
\[
  \frac{dQ}{dt}
  =
  -2\nu\|\nabla u\|_{L^2}^2
  -2\nu\|\Delta u\|_{L^2}^2
\]
has lost the nonlinear enstrophy source.
The honest inequality contains
\[
  +\,C\|\nabla u\|_{L^2}^3 .
\]
Calling \(Q(t)\) monotone at this point assumes the theorem that the proof still
has to supply.

What survived from Body II is the source-wall question:
where does the nonlinear source go?

\subsection{Body III: Non-Sobolev Compactness}

The compactness body tried to make the first two bodies produce an exact
Navier--Stokes object.
The source norm was
\[
  \|u\|_{H_{\mathrm{ns}}}
  =
  \|u\|_{L^2}
  +\|\nabla u\|_{L^2}
  +\sup_{k>k_0}\frac{E(k)}{k^\beta}.
\]
The actual compactness statement is a low/high split with explicit hypotheses.
Write
\[
  u_n=u_n^{(<K)}+u_n^{(\ge K)} .
\]
The low-frequency part has only finitely many Fourier modes and is compact once
a time-derivative bound in a negative Sobolev space is supplied.  The
high-frequency part is small only under a uniform tail condition, for example
\[
  \lim_{K\to\infty}
  \sup_n\int_0^T
  \int_{k\ge K} k^{2s}E_n(k,t)\,dk\,dt
  =0
\]
at the derivative level \(s\) being passed.  At \(s=0\), this gives
\[
  \|u_n^{(\ge K)}\|_{L^2}^2
  \le C_{LP}
  \int_{k\ge K}E_n(k,t)\,dk .
\]
Pointwise power-law decay \(E_n(k,t)\le Ck^{-\gamma}\) is only one sufficient
extra profile assumption.  The proof itself needs the displayed uniform tail
smallness, matched to the norm and time window used in the nonlinear passage.

The compactness route is real.
It tells us how a controlled approximation family can produce an exact
same-surface object:
\[
  u_n\to u
  \quad\text{strongly in}\quad
  L^2_{\mathrm{loc}} .
\]
That strong convergence is exactly what is needed to pass the quadratic term
\[
  u_n\otimes u_n \to u\otimes u .
\]

The failure is downstream dependence.
Compactness consumes the scale barrier and the \(Q(t)\) bound.
Those estimates have to arrive before compactness is used.
The derivative-control question then passes to the exact limit.
That is the fourth body.

\subsection{Body IV: Curvature Or Gradient Recovery}

The recovered fourth body is geometric.
It rewrites the equation on a Riemannian carrier:
\[
  \partial_t u+\nabla_u u
  =
  -\operatorname{grad}_g p+\nu\Delta_g u .
\]
The model gradient inequality used by this route is
\[
  \partial_t|\nabla u|^2
  \le
  -\operatorname{Ric}(g)(\nabla u,\nabla u)
  +\nu\Delta_g|\nabla u|^2
  +C_g|\nabla u|^3 .
\]
If
\[
  \operatorname{Ric}(g)\ge\kappa>0,
\]
one gets a damping term in the integral inequality:
\[
  \frac{d}{dt}\int_M|\nabla u|^2\,dV_g
  \le
  -\kappa\int_M|\nabla u|^2\,dV_g
  +C\int_M|\nabla u|^3\,dV_g .
\]

The exact mechanism is clear inside the auxiliary geometry: the positive Ricci
term supplies a geometric drain on derivative growth.
That is exactly the kind of fourth bridge the circuit wanted.

The failure is source fidelity to the Clay surface.
On the flat torus the Ricci damping term is zero.
A geometric theorem on a positive-Ricci carrier lives on a different surface
from a periodic flat-torus gradient theorem.
The fourth bridge so had to be either pulled back to the chosen
classical surface or replaced by a new same-surface theorem.

The modern replacement candidate is Lagrangian deformation geometry:
\[
  \partial_t v + A^\top\nabla_a q
  =
  \nu\,\operatorname{div}_a(G\nabla_a v),
  \qquad
  G=AA^\top .
\]
In those variables, transport is built into the geometry and viscosity
becomes heat on a moving coefficient field.
That is the strongest descendant of the old curvature idea.
It remains a theorem target until the same-surface drain law is proved.

\subsection{The First Full Failure Of The Circuit}

The intended positive loop was
\[
  \begin{aligned}
  &\text{tail control}
  \Longrightarrow Q(t)\text{ control}
  \Longrightarrow \text{compact exact object}\\
  &\Longrightarrow
  \text{gradient recovery}
  \Longrightarrow
  \text{tail recertification}.
  \end{aligned}
\]
The actual recovered state was
\[
  \begin{aligned}
  &\text{tail control demands nonlinear transfer control},\\
  &Q(t)\text{ demands payment for the enstrophy source},\\
  &\text{compactness consumes those controls},\\
  &\text{the curvature bridge lives on a different geometric surface},\\
  &\text{tail recertification remains open}.
  \end{aligned}
\]
That is the point where defects, tails, and remainders start to multiply.
The later work should be read as attempts to make the unpaid pieces visible
enough to pay.

\subsection{Modernized Four-Body Packet}

The modern route kept the four jobs and changed the readout.
The packet became
\[
  \mathcal P_{\mathrm{mod}}(T)
  =
  (S_{\mathrm{tail}},Q_{\mathrm{drain}},C_{\mathrm{exact}},G_{\mathrm{geo}};
  \mathcal E_{\mathrm{live}}^{(n)}).
\]
The desired relay was
\[
  \bigl(\mathcal P_{\mathrm{mod}}(T),
  \mathcal E_{\mathrm{live}}^{(n)}(T)\bigr)
  \Longrightarrow
  \bigl(\mathcal P_{\mathrm{mod}}(T^+),
  \mathcal E_{\mathrm{live}}^{(n+1)}(T^+)\bigr).
\]
The live survivor appears in three faces:
\[
  \text{tail leakage}
  \quad\leftrightarrow\quad
  \text{tower flux/debt}
  \quad\leftrightarrow\quad
  \text{deformation-geometry commutator}.
\]
This was progress because it stopped treating every new vocabulary as a new
route.
It said that the dyadic packet, the derivative tower, and the deformation
geometry were likely reading the same survivor.

The failure of the modern packet is also clear.
It redescribes the survivor more accurately.
The contraction theorem remains open.
The missing theorem is a same-surface drain or dissipative-margin law strong enough to
return from Body IV to a stronger Body I tail theorem.

\subsection{Two Survivor Families}

After the modern four-body pass, the written analysis had two terminal survivor families.
Family A is the lower-prefix / affine-quotient family.
Its objects are scale-memory, Volterra primitives, persistent weighted quotient
modes, and shell low-frequency strain.
Family B is the upper-tail / gap-kernel family.
Its objects are lifted-band leakage, gap-kernel flux, upper boundary packets,
and kernelized tail dissipation.

The old sentence, ``control the tail,'' had become too coarse.
For Family A, the difficult object is a one-sided memory of lower scales acting
on an active shell:
\[
  \sum_{k\le j-C} E_k(t)\,2^{-j}D_j(t)^2 .
\]
For Family B, the difficult object is an upper-tail leakage packet with a
gap-kernel gain:
\[
  \sum_{j\ge N}\sum_{N+M<k<j-4}
  2^{-(j-k)}\,2^{3k/2}\|\Delta_k u\|_2
  \|\nabla\Delta_j u\|_2\|\Delta_j u\|_2 .
\]
Both are born from the same four-body payment problem.
They fail in different ways.

The bridge obligations were:
\[
  B1:\ \text{put the lower-prefix family into a coercive drain law},
\]
\[
  B2:\ \text{put the upper-tail family into a kernelized leakage barrier},
\]
\[
  B3:\ \text{prove the two families are the same terminal survivor or make
  one pay the other}.
\]
The written analysis reached reductions in both directions.
The theorem that collapses Family A and Family B into one closed payment
remains open.
This split is why later sections keep lower-prefix memory, affine moments,
upper-tail packets, and source reserve as separate objects until a Part/Field landing
is actually derived.

\subsection{Kernelized Lifted Band}

The lifted-band route attacks the exact open shell geometry
\[
  N+M<k<j-4 .
\]
The natural scale-side carrier reached by the shell algebra is
\[
  \mathfrak K_N^{\mathrm{scale}}[u^{(n)}](t)
  =
  \sum_{j\ge N}\sum_{N+M<k<j-4}
  2^{-(j-k)}2^{3k/2}\|\Delta_k u^{(n)}\|_2\,
  \|\nabla\Delta_j u^{(n)}\|_2\,
  \|\Delta_j u^{(n)}\|_2 .
\]
The factor
\[
  2^{-(j-k)}
\]
is a real gain.
The gap kernel turns the lifted band into a separated resonance packet.

Try to integrate this carrier using Young's inequality.
Set
\[
  \lambda_j(t)
  =
  \sum_{N+M<k<j-4}
  2^{-(j-k)}2^{3k/2}\|\Delta_k u(t)\|_2 .
\]
Then
\[
  \lambda_j\|\nabla\Delta_j u\|_2\|\Delta_j u\|_2
  \le
  \varepsilon\|\nabla\Delta_j u\|_2^2
  +
  C_\varepsilon\lambda_j^2\|\Delta_j u\|_2^2 .
\]
The first term is placed under the dissipative margin.
The second term becomes the test.
By Cauchy-Schwarz in \(k\),
\[
  \lambda_j^2
  \le C_\lambda
  \sum_{N+M<k<j-4}
  2^{-(j-k)}2^{3k}\|\Delta_k u\|_2^2 .
\]
Thus the remainder is bounded by
\[
  \sum_{k\ge N+M}
  2^{3k}\|\Delta_k u\|_2^2
  \sum_{j>k+4}2^{-(j-k)}\|\Delta_j u\|_2^2 .
\]
The current energy surface controls
\[
  \int_0^T\sum_k 2^{2k}\|\Delta_k u(t)\|_2^2\,dt .
\]
The lifted-band remainder contains
\[
  2^{3k}\|\Delta_k u\|_2^2
  =
  2^k\cdot 2^{2k}\|\Delta_k u\|_2^2 .
\]
The extra factor \(2^k\) is unpaid by the present Body I inputs.

This route made a real reduction and still stopped.
The off-diagonal kernel is installed at the algebra level.
The missing theorem is an integrated \(k\)-side payment of the displayed
uncontrolled factor \(2^k\), for example a theorem that supplies that residual
square factor on the intermediate shell or an equivalent coupled weighted tail
theorem.

\subsection{Shell Gain And The Half-Derivative Target}

The shell-gain branch tried to pay the missing \(2^k\) factor by proving a
normalized low-frequency strain barrier.
Let
\[
  A_j(t)=\sum_{k\le j}2^{5k/2}\|\Delta_k\omega(t)\|_2,
  \qquad
  B_j(t)=\frac{A_j(t)}{\nu 2^{2j}} .
\]
The required active-window statement is that, for a chosen small \(\varepsilon>0\), the
active shell sees
\[
  B_j(t)\le\varepsilon
\]
on the windows where the high-frequency packet is trying to grow.
That would convert the leftover lower-shell factor into a small coefficient
against viscous loss.

The finite-shell reduction has one concrete step.
Fix \(j\).
Low shells \(k\le j-C\) have a separation gain, and the high shell has viscous
time \(2^{-2j}\).
On one shell, the desired payment is the inequality
\[
  \int_I A_j(t)\,D_j(t)\,dt
  \le
  \varepsilon\nu\int_I 2^{2j}D_j(t)\,dt
  +\mathrm{lower\ order}.
\]
The proof closes under the exact active-window condition
\(A_j(t)\le\varepsilon\nu 2^{2j}\).

The available data gives
\[
  \int_I \sum_k2^{2k}\|\Delta_k\omega(t)\|_2^2\,dt<\infty ,
\]
which is an \(L_t^1\) square budget.
The active-window bound
\[
  \sum_{k\le j}2^{5k/2}\|\Delta_k\omega(t)\|_2
  \le
  \varepsilon\nu 2^{2j}
\]
remains open.
The half-derivative target is exactly the missing improvement between those
two statements.

This branch survives as a precise demand:
\[
  \text{produce dynamic suppression of low-frequency strain at the viscous
  scale}.
\]
The repaired shell note supplies scheduled periodic shell/tail suppression.
The all-threshold endpoint half-frequency gain remains open.

\subsection{Cumulative Tail Stress And LPAS}

The far-corona route starts in continuum scale variables.
It isolates a coupled two-scale density
\[
  \mathcal A_{\ell,r}(t)
  =
  \chi_{\mathrm{mid}}(\log\ell)
  \int_{\mathbb T^3}
  \widetilde\Sigma_r(x,t)\,
  \tau_\ell^{H^1}[u](x,t)\,dx,
  \qquad
  r\ge c_1\ell .
\]
The far-corona packet is bounded by the kernel
\[
  \Pi_N^{\mathrm{mid,far},L}(t)
  \le C_L
  \int_0^{\ell_N}\int_{r\ge c_1\ell}
  \frac{\ell}{r}\mathcal A_{\ell,r}(t)
  \frac{dr}{r}\frac{d\ell}{\ell}.
\]
The exact continuum debt is to close cumulative-tail stress strongly enough to
make that integral small.

The direct inherited Euclidean reduction lands on a dyadic shadow.
Define
\[
  D_j(t)=2^{4j}\|\Delta_j u(t)\|_2^2,
  \qquad
  E_j(t)=2^{2j}\|\Delta_j u(t)\|_2^2,
\]
and
\[
  \mathcal P_j^\downarrow(t)
  =
  \sum_{k\le j-C_1}E_k(t).
\]
The lower-prefix active-square theorem asks for
\[
  \mathrm{LPAS}_N
  :=
  \int_0^T\sum_{j\ge N}
  \mathcal P_j^\downarrow(t)\,2^{-j}D_j(t)^2\,dt
  \le
  \varepsilon\nu\int_0^T D_N(t)\,dt+C2^{-2\delta N}.
\]

The proof attempt has three natural failures.
First, the installed upper-tail queues have the opposite orientation.
Expanding the lower prefix gives
\[
  \mathrm{LPAS}_N
  =
  \int_0^T\sum_{j\ge N}\sum_{k\le j-C_1}
  E_k(t)\,2^{-j}D_j(t)^2\,dt .
\]
This is lower-prefix energy multiplied by an active-shell square.
The existing high-tail queues control shells above the active index.
They leave the lower-prefix direction unpaid.

Second, the crude bound
\[
  \mathcal P_j^\downarrow(t)\le\|\nabla u(t)\|_2^2
\]
leaves
\[
  \int_0^T
  \|\nabla u(t)\|_2^2
  \sum_{j\ge N}2^{-j}D_j(t)^2\,dt .
\]
The energy law gives \(\|\nabla u\|_2^2\in L_t^1\).
It leaves the active square factor uncontrolled.

Third, pointwise coefficient-margin control would need
\[
  \mathcal P_j^\downarrow(t)\,2^{-j}D_j(t)
  \le
  c\nu
\]
on active windows.
That is an active-shell amplitude theorem.
Energy and dissipation alone leave that theorem open.

The route so stops at one explicit object:
\[
  \mathcal P_j^\downarrow\,2^{-j}D_j^2 .
\]
The continuum cumulative-tail route, the dyadic LPAS route, and the active
square route are the same payment problem seen at different resolutions.

\subsection{Volterra And Lower-Triangular Memory}

The Volterra route notices that the lower-prefix object is ordered in scale.
Set \(s=\log r\), work on a middle-band interval \(I=[s_0,s_1]\), and define
\[
  (\mathcal V_\kappa f)(s)
  =
  \int_{s_0}^{s-\kappa} f(\sigma)\,d\sigma .
\]
The far-corona primitive can be written as
\[
  F(s,t)=(\mathcal V_\kappa\widetilde Z)(s,t),
\]
and the weighted stress as
\[
  \mathrm{FCTS}
  =
  \int_0^T
  \|M_w^{1/2}\mathcal V_\kappa \widetilde Z(\cdot,t)\|_{L_s^2(I;L_x^2)}^2\,dt .
\]
The quadratic form is
\[
  \mathcal Q_w[f]
  =
  \|M_w^{1/2}\mathcal V_\kappa f\|_{L_s^2(I;L_x^2)}^2
  =
  \langle f,\mathcal V_\kappa^\ast M_w\mathcal V_\kappa f\rangle .
\]
Its kernel is
\[
  \Gamma_w(\sigma,\tau)
  =
  \int_{\max\{\sigma,\tau\}+\kappa}^{s_1}w(s)\,ds .
\]

This operator view reveals the quotient that crude Hardy estimates erase.
In the weighted space \(L^2(I,w\,ds)\), the exact-in-scale subspace is
\[
  \mathcal D_0=\{\partial_s S:S(s_0)=S(s_1)=0\}.
\]
The weighted orthogonal complement is represented by
\[
  m_w(s,t)
  =
  \frac{|I|w(s,t)^{-1}}{\int_I w(\sigma,t)^{-1}\,d\sigma}.
\]
The projector
\[
  (P_w f)(s,t)
  =
  m_w(s,t)\frac1{|I|}\int_I f(\sigma,t)\,d\sigma
\]
splits the lower-triangular output:
\[
  \mathcal V_\kappa
  =
  P_w\mathcal V_\kappa+(I-P_w)\mathcal V_\kappa .
\]
The persistent term is
\[
  \|M_w^{1/2}P_w\mathcal V_\kappa f\|^2
  =
  \frac{|I|^2}{\int_I w^{-1}\,ds}
  \left\|
    \frac1{|I|}\int_I(\mathcal V_\kappa f)(s)\,ds
  \right\|_{L_x^2}^2 .
\]

This was real progress.
The full lower prefix was reduced to one weighted quotient mode plus an
exact-in-scale remainder.

The route still failed quantitatively.
The weight \(w\) can degenerate.
No inherited estimate controls \(w^{-1}\).
The exact-in-scale remainder has no new same-surface coercive bound.
Thus the route identifies the survivor:
\[
  \text{weighted quotient mode / affine moment pair }(M_0,M_1),
\]
and leaves its vanishing as a new theorem.

\subsection{Signed Descent And The Limits Of Cancellation}

The signed descent route tried to use cancellation before taking absolute
values.
It replaced dyadic shells by a continuous Littlewood--Paley scale:
\[
  P_{\le \ell}=\varphi(\ell D),
  \qquad
  Q_\ell=-\ell\partial_\ell P_{\le \ell}.
\]
The signed lifted commutator density was
\[
  \mathfrak c_\ell(x,t)
  =
  \ell^{-2}
  u_\ell(x,t)
  [Q_\ell,b_\ell^{\mathrm{meso}}\cdot\nabla]u(x,t).
\]
The key observation is that \(Q_\ell=-\ell\partial_\ell P_{\le\ell}\).
So the lifted commutator appears as a derivative in scale of a cumulative
potential:
\[
  \ell\partial_\ell \mathcal A_\ell
  =
  \mathfrak c_\ell
  +\mathfrak c_\ell^{\mathrm{transpose}}
  +\mathfrak s_\ell^{\mathrm{edge}}
  +\mathfrak w_\ell .
\]
After the edge and weight terms are routed, the desired signed form is
\[
  \mathfrak c_\ell
  =
  (\ell\partial_\ell)\Psi_\ell
  +\nabla_x\cdot J_\ell
  +\partial_t R_\ell
  +\mathcal S_\ell^{\mathrm{spill}}
  +\mathcal E_\ell .
\]
Integrating in \(x,t,\ell\) would telescope the scale derivative and leave only
boundary, spill, and residual terms.

That route explains the correct order:
\[
  \text{signed density}
  \Longrightarrow
  \text{signed telescoping}
  \Longrightarrow
  \text{residual majorant}.
\]
The positive majorant should enter only after the signed mechanism has done its
work.

The direct signed weighted theorem still fails.
The dyadic lifted remainder has principal form
\[
  \mathcal R_{j,k,\ell}^{\mathrm{lift}}
  =
  2^{2j}
  \langle[\Delta_j,a_k\cdot\nabla]\Delta_\ell u,\Delta_j u\rangle,
  \qquad
  a_k=\Delta_k u .
\]
For \(k<j\) and \(|\ell-j|\le C_{\mathrm{col}}\),
\[
  [\Delta_j,a_k\cdot\nabla]\Delta_\ell u
  =
  2^{-j}\mathcal C_j(\nabla a_k,\nabla\Delta_\ell u)
  +\text{collar terms}.
\]
The leading term is a strain form.
It has no fixed sign.
Active wave packets can align with expanding eigendirections of the symmetric
part of \(\nabla a_k\).
The one-sided weights orient the shell sum.
Eigendirection cancellation remains open.

Thus the signed branch stopped at a specific missing theorem:
active eigendirection decorrelation, signed measure cancellation, or a
structural nonpositive divergence identity.
Taking absolute values returns to the active-square / source-reserve wall.

\subsection{Persistent Quotient Kill Attempts}

Once the Volterra route exposed the affine quotient, three kill routes were
tried.
Route A pushed the quotient back into cumulative-tail stress geometry.
That collapsed the problem to the same CTS/LPAS payment:
\[
  \mathcal Q_w[f]
  \leadsto
  \int_0^T\sum_j\mathcal P_j^\downarrow\,2^{-j}D_j^2\,dt .
\]
Route C used far-corona tail structure to remove the quotient.
That route became circular because the far-corona theorem was one of the
statements the quotient was supposed to help prove.

Route B was the serious reduction.
It isolated two affine moments,
\[
  M_0(t)=\int_I f(s,t)\,ds,
  \qquad
  M_1(t)=\int_I s\,f(s,t)\,ds .
\]
If both moments vanish, the weighted quotient loses its persistent mode.
The signed descent identity can control scale derivatives of cumulative
potentials, so it looked like the correct tool for killing \(M_0\) and \(M_1\).

The proof reaches the boundary of what the signed identity provides.
The identity controls a scale derivative:
\[
  \ell\partial_\ell\Psi_\ell
  =
  \mathfrak c_\ell-\nabla_x\cdot J_\ell-\partial_tR_\ell
  -\mathcal S_\ell^{\mathrm{spill}}-\mathcal E_\ell .
\]
The affine moments require one more integrated scale payment.
The boundary terms at the ends of the scale interval, together with spill and
edge terms, can carry the same affine mass the proof is trying to erase.
Thus signed descent reduces the quotient to a smaller pair of moments, then
stops one scale derivative short of killing them.

\subsection{Projected Flow, TPS, And SG.4}

The projected-flow sidecar and the two-point strain packet both tried to put
the defect into a parabolic equation.
The projected-flow order was:
\[
  D.1\ \text{well-posedness},\qquad
  D.2\ \text{coercive energy},\qquad
  D.3\ \text{commutative shadow},
\]
\[
  D.4\ \text{shadow continuation},\qquad
  D.5\ \text{Clay integration}.
\]
The projected-flow energy identity has the shape
\[
  \frac{d}{dt}E_D(X)
  +\text{carrier dissipation}
  +\text{commutator-tier dissipation}
  =
  N_1(X)+N_2(X),
\]
with the required estimate
\[
  |N_1|+|N_2|
  \le
  \varepsilon\,\text{dissipation}+C_\varepsilon\Psi(E_D).
\]
The route requires a projected noncommutative flow with a coercive energy and a
commutative shadow: prove the projected theorem, shadow it back to the ordinary
equation, and inherit continuation.
The open points are \(D.1\), the \(D.2\) bilinear hinge, and global
\(H_D^s\) continuation.

The two-point strain route defines
\[
  W_J=K_J-\overline K_J
\]
and gives \(W_J\) a parabolic equation with drift, diffusion, and remainders.

The structural match was real:
\[
  \text{drift}+\text{diffusion}+\text{remainder}
  \Longrightarrow
  \text{defect estimate}.
\]
The split came from the readout.
SG.4 asked for a one-sided lower envelope on selected directional strain:
\[
  \frac{d}{dt}\log d_{ab}^{\mathrm{coarse}}(t)
  =
  s_J(a,b,t),
  \qquad
  s_J(a,b,t)\ge\lambda_J(t)-\varepsilon_J(a,b,t).
\]
The exact pair-distance identity is strong:
\[
  \frac{d}{dt}\log |X(a,t)-X(b,t)|
  =
  e_{ab}(t)\cdot S(X(a,t),t)e_{ab}(t)
  +\mathrm{commutator/error}.
\]
The SG.4 production theorem would need to select pairs for which the
right-hand side has the lower bound \(\lambda_J-\varepsilon_J\) often enough
to renew the barrier.
A norm can control size and still lose sign, direction, selector membership,
and one-sided growth.
The selector-observability audit gives the no-go point:
the same signless \(W_J\) data can represent different expanding-direction
deficits.
D.2/TPS support so leaves SG.4 open.

What survived is the defect-object viewpoint.
The missing theorem is selector-good production from the one-sided strain
budget.

\subsection{The Source-Wall Deletion Cycle}

The \(Q(t)\) body exposed the nonlinear enstrophy source.
The next long phase tried to delete or pay that source directly.
The pattern was stable:
\[
  \text{find source pulse or reserve}
  \Longrightarrow
  \text{try to charge it}
  \Longrightarrow
  \text{find a sharper survivor}.
\]

The active source-wall root was expressed as
\[
  \text{ScaleCriticalTreeCarleson.A}
  \quad\text{or}\quad
  \text{ZenoBoundedClass}_{B}\text{.A}
  +\text{ZenoResidueLiouville}_{B}\text{.A}.
\]
Equivalently, the route wanted native positive source control:
\[
  \int_I
  \sum_{P\in\mathcal F_N}
  \bigl(\mathfrak S^{\mathrm{active}}_{N,\mathrm{loc}}(P,t)\bigr)_+\,dt
  \le
  o_N(1)+\mathrm{Loss}_{\mathrm{legal}} .
\]
Taking absolute values returned to the scale-critical tree reserve.
Passing to compactness returned to the Zeno residue route.
Thus the source wall kept reappearing under several names.

\subsection{BASAC And The Terminal Atom}

The BASAC branch produced a terminal tangent class with routed legal exits,
zero ASAC defect, no finite donor tree left, and a same-fluid terminal source
atom.
The model time density is the key:
\[
  g_m(s)=m\,\mathbf 1_{(-1/m,0]}(s).
\]
It has bounded \(L_s^1\) mass:
\[
  \int_{-1}^0 g_m(s)\,ds=1,
\]
and it converges weakly to a terminal atom at \(s=0\).
For every \(p>1\),
\[
  \int_{-1}^0 g_m(s)^p\,ds
  =
  m^{p-1}
  \to\infty .
\]
For every superlinear Orlicz function \(\Phi\),
\[
  \int_{-1}^0\Phi(g_m(s))\,ds
  =
  \frac1m\Phi(m)
  \to\infty
\]
along a subsequence.

This calculation explains the failure.
The BASAC bookkeeping gives \(L^1_t\) source mass and a terminal same-fluid
ledger.
Temporal uniform integrability, reverse Holder residence, Morrey decay, and
other anti-atom moduli remain open.

A conditional theorem exists.
If one produces a subclass \(B_{\mathrm{ASAC}}^{UI}(\Phi)\) with
\[
  \sup_m\int_{-1}^0\Phi(g_m^R(s))\,ds<\infty ,
\]
then de la Vallee-Poussin gives uniform integrability and the terminal source
atom vanishes.
The production theorem
\[
  B_{\mathrm{ASAC}}\Longrightarrow B_{\mathrm{ASAC}}^{UI}(\Phi)
\]
is the missing step.

The later CM use is different.
The terminal atom becomes a tested obstruction.
After the CM test admits the terminal object, a zero-radius source residue is
sorted first through the Part/Field side.
That CM use records what the surviving object means once it is offered as a
finite terminal obstruction.
The positive BASAC deletion theorem remains a separate supplier theorem.

\subsection{Pressure Sustain And Pure Pressure Residue}

The pressure route tried to tether pressure-sustain residue to native source
residue.
The desired bridge was
\[
  \text{PressureSourceAC.A}.
\]
At the limit, the pressure residue splits as
\[
  \mu_*^{\mathrm{press\text{-}sus}}
  =
  f\Lambda_*^{\mathrm{src}}
  +
  \mu_*^{\mathrm{press}\perp\mathrm{src}},
  \qquad
  \Lambda_*^{\mathrm{src}}
  =
  \nu_*^{\mathrm{src}}+\nu_*^{\mathrm{legal,src}} .
\]
The source-carried part can join the source-wall ledger.
The singular part
\[
  \mu_*^{\mathrm{press}\perp\mathrm{src}}
\]
is orthogonal to the native source/legal source carrier.

The failure is that pressure Poisson structure gives a quadratic-gradient
pressure carrier distinct from a native same-fluid source-current carrier.
Calderon--Zygmund relations act spatially and tensorially.
The transported source ancestry needed by the native source-wall argument
remains open.

The route ends at a clean fork:
\[
  \text{PressureSourceAC.A}
  \quad\text{or}\quad
  \text{PurePressureSustainResidueLiouville.A}.
\]
Without one of those the pressure survivor is only a pressure-side diagnostic.
It can later inform a Field-face analysis after the pressure-specific exits are
isolated.
Native source deletion requires the tether or Liouville theorem.

\subsection{Angular Saturation}

The angular pressure attempt used the trace and angular cancellation in the
pressure Hessian.
For a selected direction \(n_P\), the positive lobe is
\[
  [-n_P\cdot\nabla^2p\,n_P]_+ .
\]
Trace and Calderon--Zygmund cancellation say that the full pressure Hessian has
balanced directional content.
The terminal ledger, however, keeps one selected positive part with terminal
weights and filters.

The direct proof fails because negative angular partners can be diffuse,
subthreshold, mismatched in weights, or routed outside the retained terminal
family.
Full tensor cancellation can coexist with one selected positive pressure bill.

The surviving pressure theorem is a partner theorem:
\[
  \begin{aligned}
  &\text{TerminalPressurePartnerSaturation.A}
  +\text{CZPressurePairWeightDefectCharge.A}\\
  &+\text{TerminalPressureCollarLegal.A}
  \Longrightarrow
  \text{TerminalPressureTraceAngularSaturation.A}.
  \end{aligned}
\]
The missing branch is again a pressure lobe/source tether or pure pressure
Liouville theorem.

\subsection{Transported Boundary And No Incoming Flux}

Transported cylinders are the right geometry for same-fluid ancestry.
They remove artificial fixed-frame boundary flux.
The required theorem was
\[
  \mathrm{Flux}^{\mathrm{src},+}_N
  (\partial\mathcal C_R\times[T-\varepsilon,T])
  \to0
\]
as the terminal layer shrinks.

Boundary control alone leaves an interior time-face atom available.
The model obstruction is
\[
  F_N^{\mathrm{src},+}(t)
  =
  \tau_N^{-1}\mathbf 1_{[T-\tau_N,T]}(t)
\]
supported inside the transported cylinder.
It produces terminal source mass while boundary inflow remains small.

A conditional local-energy bridge exists.
On a same-fluid transported cylinder,
\[
  \Delta E_{\mathrm{loc}}^{\mathrm{term}}
  +D_{\mathrm{loc}}^{\mathrm{term}}
  =
  \mathrm{Source}_{\mathrm{native}}^{\mathrm{term}}
  +\mathrm{Flux}_{\mathrm{in}}^{\mathrm{term}}
  +\mathrm{Legal}^{\mathrm{term}}
  +\mathrm{ASAC}^{\mathrm{term}}
  +\mathrm{Donor}^{\mathrm{term}} .
\]
If
\[
  \mathrm{Flux}_{\mathrm{in}}^{\mathrm{term}}=0
\]
and the positive terminal local-energy jump vanishes, then a positive terminal
source atom has no balancing channel and is excluded.

The missing production theorem is the full transported no-incoming flux plus
interior production exclusion.
Transported geometry helps boundary bookkeeping.
Time concentration still needs a separate spreading, trace, or
anti-concentration theorem.

\subsection{Quantized Parent And Zeno Ancestry}

The ancestry route tried to make infinite same-fluid refill spend a quantized
resource.
A useful theorem would say that each legal terminal refill edge spends at least
\(\eta>0\) of some monotone quantity.
Then infinitely many refills would contradict a finite original-data budget.

Each candidate resource failed:
\[
  r_k\downarrow0
\]
orders the scales and gives no uniform integer decrement.
Heat-time depth allows
\[
  \sum_k r_k^2<\infty ,
\]
which is exactly compatible with infinitely many terminal scales.
Raw energy and dissipation can split over infinitely many descendants.
Zero ASAC defect has already removed that payment channel.
Same-fluid labels preserve ancestry and leave the decreasing alphabet absent.
Finite donor balance telescopes only finite truncations and leaves the terminal
leaves.

The direct quantization theorem so fails from the installed inputs.
The reduced alternatives are:
\[
  \begin{gathered}
  \text{ScaleNormalizedBranchEntropy.A},\qquad
  \text{QuantizedParentCharge.A},\\
  \text{MinimalBadAncestryCompactness.A with rigidity}.
  \end{gathered}
\]
The route ends at a terminal Zeno residue unless a genuine rigidity or Liouville
class is produced.

\subsection{Local Energy And Elliptic Time Smearing}

The local-energy branch tried to use the local energy trace to forbid terminal
source atoms.
The signed local energy inequality gives compactness and admissible terminal
defect measures.
It leaves the selected positive native source component uncontrolled.
The positive part can sit in a shrinking time slab while the signed balance
remains legal.

The elliptic time-smearing branch tried to use pressure ellipticity or
Calderon--Zygmund spreading.
The obstruction is simple.
Spatial elliptic operators act on each time slice:
\[
  \rho(s,y)\longmapsto T_y\rho(s,y).
\]
If
\[
  \rho_m(s,y)=m\,\mathbf 1_{(-1/m,0]}(s)\rho(y),
\]
then the spatial operator changes \(\rho(y)\) and preserves the time support
\((-1/m,0]\).
The terminal atom in time survives.

Thus local energy and elliptic pressure support source-wall analysis.
The native terminal atom remains unless a separate trace, flux, or
anti-concentration theorem is supplied.

\subsection{All Reduced Source-Wall Routes}

The all-routes source-wall pass left six forward supplier candidates:
\[
  \text{UniformTemporalSourceIntegrability}_{p,B_{\mathrm{ASAC}}}\text{.A},
\]
\[
  \text{LocalPositiveSourceCarleson.A / PositiveRemainderDepletion.A},
\]
\[
  \text{QuantizedParentCharge.A / MinimalBadAncestryCompactness.A with rigidity},
\]
\[
  \text{PressureSourceAC.A / LocalizedLeraySourceTether}_{ind}\text{.A},
\]
\[
  \text{PurePressureSustainResidueLiouville.A},
  \qquad
  \text{TerminalPressurePureStrainLegal.A}.
\]
No one of these was installed unconditionally from original smooth data.

That is the source-wall failure list in its reduced form.
It says the forward deletion program has stopped at real mathematical doors,
with explicit equations and missing theorems.
The shared residue is the terminal zero-thickness native source concentration,
plus pressure-side variants that need their own tether or Liouville theorem.

\subsection{Source Reserve Birth}

The May source-reserve burst sharpened the source wall.
Choose the first terminal same-fluid window \(W\) where donor-refill reserve
surplus is born, and write
\[
  R_N(W)
  =
  \int_W\sum_{k>N}2^k
  \left(\sum_{\ell>k+4}D_\ell(t)\right)^2\,dt .
\]
The desired charge inequality was
\[
  \bigl[
    R_N(W)
    -(1-\delta)R_N(\mathrm{Past}(W))
    -\mathrm{Loss}_{\mathrm{legal}}(W)
  \bigr]_+
  \le
  \mathrm{Charge}^{\mathrm{native}}_N(W)+o_N(1).
\]
The native charge had to land on the same retained terminal ledger:
legal packing payment, participation failure, licensed Field source charge, or
terminal Zeno residue already sorted in the witness tree.

The proof chose \(W\) minimally and decomposed the reserve into
\[
  \text{inherited past reserve}
  +\text{first-created native reserve}
  +\text{child residual reserve}
  +\text{legal loss reserve}.
\]
Minimality contracts the inherited term.
Existing residual and legal ledgers route the child and legal terms.
The survivor is
\[
  \text{first-created native reserve}\ge c\eta .
\]

The direct proof stops here.
The first-created positive-scale native donor-square reserve remains unpaid by
the installed ledgers.
Zero-radius donor ancestry is already a carrier-window support boundary.
The retained positive-scale first-created reserve remains supplier-quarantined
until a same-ledger payment theorem or a CM-test admission plus first Part/Field
failure is proved.

\subsection{Zero-Moment Signed Pair}

The zero-moment attempt sharpened the source-reserve survivor again.
It showed that zero first moment can coexist with mass bounded below by a fixed
positive constant.
On the first terminal window,
\[
  R_{\mathrm{first}}(W)\ge\eta,
  \qquad
  M_{\mathrm{first}}(W)=0 .
\]
Split the source carrier as
\[
  J=J_+-J_-,
  \qquad
  \int_WJ_+=\int_WJ_-,
  \qquad
  c_J\le \int_W(J_++J_-)\le C_J .
\]
The square reserve can see a pair with
\(c_J\le \int_W(J_++J_-)\le C_J\) while the signed first moment
vanishes.

Most cases route cleanly.
A partner on the wrong carrier breaks Part or Field.
A partner that survives only through terminally shrinking windows becomes Zeno.
The hard case is
\[
  J_+\text{ and }J_-
  \text{ survive at positive scale on the same retained ledger,}
\]
with square reserve bounded below by \(\eta>0\) and zero first moment.

The failed implication says that one packet must produce visible active height flux.
That implication is too fast.
The packets can exchange through the same retained ledger while the first
moment cancels on every heat subwindow.
The missing theorem is a same-ledger signed-pair no-free-sink theorem, or an
equivalent two-tower donor depletion theorem.

\subsection{Public \(L^3\), Duhamel, And Native Source}

The public critical \(L^3\) branch has to be read more carefully than the
source-wall deletion branches.
The useful bridge is a local CM-facing translator for terminal \(L^3\) Duhamel
mass on the same witness ledger.

Start with terminal critical concentration on a ledger
\[
  \mathcal W=(\Owork,Q,N,r,\Phi,\mathcal T)
\]
and a selected concentration
\[
  \ell_m^{-2}\int_{Q_m}|u_m|^3\ge c_0 .
\]
On the terminal scale, Duhamel gives
\[
  u_m(t)
  =
  e^{(t-t_0)\nu\Delta}u_m(t_0)
  +
  \int_{t_0}^{t}
  e^{(t-s)\nu\Delta}
  P\nabla\cdot(u_m\otimes u_m)(s)\,ds .
\]

The heat side is checked by same-ledger heat ancestry.
The nonlinear side is localized by the adjoint heat/Leray response and then
captured by the paraproduct native-source argument.
The installed branch-level conclusion is
\[
  \text{terminal }L^3\text{ Duhamel response mass}
  \Longrightarrow
  \text{carrier-window support outside the exit conclusion}
  \vee
  \neg\Part_{N,Q}
  \vee
  \Legal .
\]
With retained carrier support and Part, the local public translator gives the Field-facing
critical readout.

This is a different role from the positive supplier route.
It proves the same-ledger terminal branch:
once terminal \(L^3\) Duhamel mass is already present on the same ledger, the
local branch is routed into the CM witness grammar.

\subsection{Terminal Leray Saturation}

The terminal Leray branch sits next to the \(L^3\) translator.
After legal commutator removal, the retained source work has the form
\[
  -\int (w_j\otimes\widetilde w_j):\nabla S^{\mathrm{loc}}_{<j}u .
\]
The desired saturated square is
\[
  -\int (w_j\otimes w_j):\nabla S^{\mathrm{loc}}_{<j}u
  =
  \int\langle S^{\mathrm{loc}}_{<j}w_j,w_j\rangle .
\]
Writing
\[
  \widetilde w_j=w_j+\delta_j
\]
leaves the signed-partner defect
\[
  -\int (w_j\otimes\delta_j):\nabla S^{\mathrm{loc}}_{<j}u .
\]

In the local \(L^3\) translator, the same-family signed-partner defect is
handled by the installed signed-partner no-free-defect theorem.
That closes the local terminal \(L^3\) route.

The independent pressure-Leray route remains open.
An elliptic pressure lobe still needs a transported source/donor edge theorem.
The remaining pressure-side residue is
\[
  \text{PressureLobeSourceTether.A}
  \quad\text{or}\quad
  \text{PressureSustainResidueToSourceResidue.A}.
\]

\subsection{Critical \(H^{1/2}\)}

The \(H^{1/2}\) audit is conditional.
It remains a candidate translator.
The proposed witness is a selected critical dyadic signal:
\[
  \sum_{|j-J_m|\le C}
  2^j\|P_j u(t_m)\|_{L^2(Q_m)}^2
\]
with a positive terminal lower bound or normalized divergence.

The physical split is useful.
\(L^3\) sees critical size.
\(H^{1/2}\) sees critical size together with frequency placement.
So an \(H^{1/2}\) terminal signal has two readings.

The amplitude reading routes to the installed local \(L^3\) translator.
The pure oscillatory reading keeps amplitude under control and concentrates the
defect in frequency organization.
With retained carrier support and Part, that pure oscillatory branch should be a Field
failure: neighboring towers lose coherent agreement at the selected depth.

The missing lemmas are exact:
\[
  \begin{gathered}
  \text{SameLedgerHHalfExtraction.A},\qquad
  \text{HHalfOscillationFieldExit.A},\\
  \text{HHalfAmplitudeToL3Translator.A}.
  \end{gathered}
\]
Until those are proved, \(H^{1/2}\) remains a conditional translator.

\subsection{What The Failed Attempts Have In Common}

Across the positive branches the same pattern recurs:
\[
  \begin{aligned}
  &\text{remove local or legal pieces}\\
  &\Longrightarrow
  \text{isolate an edge, collar, source, or reserve packet}\\
  &\Longrightarrow
  \text{convert the residue to a one-sided or terminal object}\\
  &\Longrightarrow
  \text{find a smaller survivor}.
  \end{aligned}
\]
The survivors have structure:
\[
  \begin{gathered}
  \text{half-derivative lower-shell deficit},\\
  \text{lower-prefix active square},\\
  \text{weighted quotient mode and affine moments},\\
  \text{same-ledger signed-pair reserve},\\
  \text{terminal zero-thickness native source atom},\\
  \text{pure pressure-sustain residue},\\
  \text{conditional critical }H^{1/2}\text{ field defect}.
  \end{gathered}
\]
Each survivor is a real mathematical stopping point for the direct positive
proof.
Each proposed payment exposed a sharper unpaid object.

\subsection{How The Class-Membership Route Uses The Failure Record}

The class-membership route uses the open positive estimates as the failure
record.
It changes the proof role of the bad terminal object.
Given an obstruction \(O\) reached by the positive program, the split is
\[
  O_{\mathrm{pass}}
  \qquad\text{and}\qquad
  O_{\mathrm{fail}}.
\]
The pass branch is the continuation branch:
\[
  O_{\mathrm{pass}}
  \Longrightarrow
  \Part_{N,Q}\wedge\Field_{N,r,Q}
  \Longrightarrow
  \Member(Q;\Owork).
\]
The fail branch must do real work.  Carrier loss remains carrier-window support
before admission; an admitted obstruction must land in Part or Field:
\[
  O_{\mathrm{fail}}
  \Longrightarrow
  \text{CM-test-admissible terminal obstruction}
  \Longrightarrow
  \neg\Part_{N,Q}
  \vee
  \forall r>0\,\neg\Field_{N,r,Q}.
\]
Only after the first Part/Field failure is derived does the class-exit conclusion
apply:
\[
  \Exit(Q;\Owork):=\neg\Member(Q;\Owork).
\]

This is why the failed attempts belong in the paper.
They identify the actual objects a finite-time breakdown claim would have to
carry.
The CM proof then asks whether that same terminal object remains in the
Navier--Stokes class or loses Part or Field.

\subsection{Branch Role Table}

\begin{center}
\begin{tabular}{L{0.27\linewidth}L{0.31\linewidth}L{0.31\linewidth}}
\textbf{Branch} & \textbf{Forward residue} & \textbf{Later CM role} \\
\hline
Four-body scale barrier &
Unproved nonlinear high-frequency transfer law &
Source of tail, lifted-band, and scale-memory witness objects \\
\(Q(t)\) / enstrophy &
Unpaid vortex-stretching source &
Ancestor of source-wall and native source-reserve branches \\
Compactness &
Exact object only after tail and energy controls are supplied &
Same-surface packet capture support \\
Curvature bridge &
Positive-Ricci geometric theorem on the wrong surface for the torus branch &
Motivation for deformation-geometry Body IV, pending same-surface drain theorem \\
Kernelized lifted band &
Residual \(2^k\) lower-shell factor &
Positive-route survivor; possible Field/tower readout after terminal admission \\
LPAS / cumulative tail &
Lower-prefix active square &
Positive-route survivor; support for source-reserve and scale-memory analysis \\
Volterra memory &
Weighted quotient mode / affine moments &
Reduced terminal survivor candidate; class exit requires Part/Field landing \\
Signed descent &
Eigendirection sign obstruction &
Returns to active-square/source-reserve wall when absolute values are taken \\
BASAC &
Terminal zero-thickness native source atom &
Zero-radius Part/Field witness support after CM-test admission \\
Pressure sustain &
Pure pressure residue orthogonal to native source &
Field diagnostic after pressure exits are isolated \\
Transported no-flux &
Interior terminal atom with small boundary flux &
Part/Field support depending on retained carrier and flux role \\
Quantized parent / Zeno &
No quantized spend; terminal ancestry can shrink indefinitely &
Zero-radius residue lands first at carrier-window support boundary unless a positive-scale carrier is retained \\
Public local \(L^3\) &
Terminal Duhamel response mass &
Installed local CM-facing witness translator \\
Critical \(H^{1/2}\) &
Same-ledger extraction and Field landing was recorded as historical pressure before the current CM finite-obstruction inventory gate passed &
Conditional translator only \\
\end{tabular}
\end{center}

\subsection{Closure Of The Failure Record}

The failed positive attempts above leave only a small number of mathematical
survivors.
Every route that tried to prove smoothness either asked for a direct estimate
closing the classical proof, or else exposed a terminal object that has to be
tested as the same Navier--Stokes solution.

The surviving objects have the same shape from different directions.
The scale route leaves unpaid high-frequency transfer.
The \(Q(t)\) route leaves the vortex-stretching source.
Compactness leaves the need for a same-surface exact object.
The curvature route leaves the need for a periodic replacement for the geometric
drain.
The later source, pressure, Zeno, \(L^3\), Leray, and critical-norm routes leave
terminal packets, residues, or readouts that must either pay one of those debts
or enter the class-membership test.

This is the point where the positive history becomes useful to the
contrapositive proof.
It asks a mathematical question about the object left by the failed route:
does the object still carry a positive same-solution packet, does that packet
still participate in the Navier--Stokes law, and does the resulting field still
have a positive-scale coherent readout?

When carrier admission is paid and the two admitted-record questions pass, the branch is the pass side.
It gives the same-solution membership branch.
When the admitted answer first fails, the branch supplies the Part/Field failure:
\[
  \neg\Part_{N,Q},
  \qquad
  \forall r>0\,\neg\Field_{N,r,Q}.
\]
Those are the only CM exit diagnostics used by the class proof, and both are
direct presentations of nonparticipation.
The record has a narrower job: it is the intake map for terminal evidence.
Each failed smoothness strategy
names the debt it could not pay, preserves the solution-surface on which that
debt lives, and hands the remaining object to the same Part/Field test
without changing the claim being tested.

Thus the failure record has a precise job.
It identifies what a finite-time breakdown claim has to carry at the terminal
time.
The main proof then tests that terminal object.
That test supplies the exhaustive negative side of Silver: non-smoothness
loses membership and full VPI participation.  Contrapositively, the original
VPI participant remains a member and is smooth.  Carrier loss remains Pack
admission support outside CM.
